\documentclass[11pt]{article}

\usepackage[a4paper,margin=27mm]{geometry}
\usepackage{amsmath,amssymb,amsthm,mathtools}
\usepackage{enumitem}
\usepackage{hyperref}
\usepackage{xcolor}
\usepackage{tikz-cd}

\hypersetup{
  colorlinks=true,
  linkcolor=blue!50!black,
  citecolor=blue!50!black,
  urlcolor=blue!50!black
}

% -----------------------------------------------------------------------------
% Basic notation
% -----------------------------------------------------------------------------
\newcommand{\USD}{\mathsf{USD}}
\newcommand{\Base}{\mathsf{Base}}
\newcommand{\RecUSD}{\mathsf{Rec}_{\USD}}
\newcommand{\St}{\mathsf{St}}
\newcommand{\UFib}{\mathcal U}
\newcommand{\DFib}{\mathcal D}
\newcommand{\Bnd}{\mathsf{Bnd}}
\newcommand{\Occ}{\mathsf{Occ}}
\newcommand{\Hand}{\mathsf{Hand}}
\newcommand{\FreeHand}{\mathsf{FreeHand}}
\newcommand{\Anch}{\mathsf{Anch}}
\newcommand{\Smart}{\mathsf{Smart}}
\newcommand{\Rel}{\mathsf{Rel}}
\newcommand{\BVis}{\mathsf{BVis}}
\newcommand{\Obs}{\mathsf{Obs}}
\newcommand{\Adopt}{\mathsf{Adopt}}
\newcommand{\NonAdopt}{\mathsf{NonAdopt}}
\newcommand{\Collapse}{\mathsf{Collapse}}
\newcommand{\IdLic}{\mathsf{IdLic}}
\newcommand{\Diff}{\mathsf{Diff}}
\newcommand{\Gap}{\mathsf{Gap}}
\newcommand{\Seq}{\mathsf{Seq}}
\newcommand{\Rev}{\mathsf{Rev}}
\newcommand{\Flip}{\mathsf{Flip}}
\newcommand{\Turn}{\mathsf{Turn}}
\newcommand{\Elim}{\mathsf{Elim}}
\newcommand{\Std}{\mathsf{Std}}
\newcommand{\Lan}{\operatorname{Lan}}
\newcommand{\Ran}{\operatorname{Ran}}
\newcommand{\colim}{\operatorname{colim}}
\newcommand{\limi}{\operatorname{lim}}
\newcommand{\op}{\mathrm{op}}
\newcommand{\xto}[1]{\xrightarrow{#1}}
\newcommand{\xot}[1]{\xleftarrow{#1}}
\newcommand{\defeq}{\mathrel{:=}}
\newcommand{\uconn}{\mathrel{\leadsto_{\USD}}}
\newcommand{\unadopt}{\mathrel{\dashrightarrow_{\USD}}}
\newcommand{\CComp}{\mathsf{CComp}}
\newcommand{\DComp}{\mathsf{DComp}}
\newcommand{\EComp}{\mathsf{EComp}}
\newcommand{\SComp}{\mathsf{SComp}}
\newcommand{\TComp}{\mathsf{TComp}}
\newcommand{\UComp}{\mathsf{UComp}}
\newcommand{\Comp}{\mathsf{Comp}}
\newcommand{\Dec}{\mathsf{Dec}}
\newcommand{\Fusion}{\mathsf{Fusion}}
\newcommand{\Fission}{\mathsf{Fission}}
\newcommand{\USDFusion}{\mathsf{USDFusion}}
\newcommand{\USDFission}{\mathsf{USDFission}}
\newcommand{\Sens}{\mathsf{Sens}}
\newcommand{\Art}{\mathcal A_{\mathrm{art}}}
\newcommand{\MagicSens}{\mathsf{MagicSens}}
\newcommand{\PlotSens}{\mathsf{PlotSens}}
\newcommand{\HallSens}{\mathsf{HallSens}}
\newcommand{\Cat}{\mathbf{Cat}}
\newcommand{\Pos}{\mathbf{Pos}}
\newcommand{\Set}{\mathbf{Set}}
\newcommand{\Prof}{\mathbf{Prof}}
\newcommand{\Arr}{\mathsf{Arr}}
\newcommand{\USDCalc}{\mathsf{USDCalc}}

\theoremstyle{definition}
\newtheorem{definition}{Definition}[section]
\newtheorem{construction}[definition]{Construction}
\newtheorem{example}[definition]{Example}
\newtheorem{remark}[definition]{Remark}
\newtheorem{assumption}[definition]{Assumption}

\theoremstyle{plain}
\newtheorem{proposition}[definition]{Proposition}
\newtheorem{theorem}[definition]{Theorem}
\newtheorem{lemma}[definition]{Lemma}
\newtheorem{corollary}[definition]{Corollary}

\title{The Up-Side-Down Functor:\
A Reference Implementation of Smart Functors with Boundary-Visible Stance}
\author{Shuhei Ihara}
\date{\today}

\begin{document}
\maketitle

\begin{abstract}
The Up-Side-Down functor is not a new equality, a new physical law, or a new
proof principle.  It is a structured functor: a pure functor equipped
with an adoptive face, a non-adoptive face, a boundary witness, and a
coherence contract.  The purpose of this paper is to give a reference
implementation of such structured functors and of their safe composition
operations.  The construction treats the adoptive and non-adoptive
regions not as two arbitrary pre-existing categories, but as fibers of a
single USD record fibration realized over both a carrier category and a
two-sided stance category.  In this
sense a USD functor is a categorical analogue of a smart pointer: after
forgetting the additional fields it is an ordinary functor, while before
forgetting it carries explicit bookkeeping for observation, non-adoption,
boundary visibility, and typed composition.  Six causeton types are
specified: compression, downward restriction, extension, sideways
stabilization, temporal turnover, and upward image.  Each causeton is
presented as a composition--decomposition pair with type conditions,
existence conditions, residue accounting, and philosophical reading.  The
conservative intent is made explicit: USD records may organize proof
search or interpretation, but USD-free conclusions must be justified in
the base theory.
\end{abstract}

\tableofcontents

% -----------------------------------------------------------------------------
\section{Scope and non-claims}
% -----------------------------------------------------------------------------

This draft responds to a common failure mode in early USD notation: terms
such as relation, boundary, observation, adoption, refusal, gap, and turn
can sound like explanatory prose rather than mathematical data.  The
present version therefore treats USD theory as a typed reference
implementation.  The reader should not be asked to infer the missing
interfaces.  They are made explicit.

\begin{remark}[What this paper does not claim]
This paper does not claim that USD notation proves a new equality, a new
inequality, a physical theorem, or a computational separation.  It does
not identify boundary-visible connection with equality.  It does not
identify non-adoption with ordinary negation.  It does not claim that the
relativistic discussion below proves anything about special or general
relativity.  The relativistic section is a test model for the vocabulary
of observers, frames, local restrictions, and temporal records.
\end{remark}

\begin{remark}[What this paper does claim]
This paper claims that one can specify a category of smart functors whose
objects are ordinary functors equipped with additional USD structure, and
that the main USD composition operations can be given as typed operations
using standard categorical vocabulary: pushouts or gluing for compression,
pullback or reindexing for downward composition, factorization for
extension, fixed points for sideways stabilization, involution plus
sequence reversal for temporal turnover, and Kan-style images for upward
composition.
\end{remark}

The intended comparison is with a smart pointer in programming.  A raw
pointer stores an address.  A smart pointer stores an address together
with ownership, lifetime, and release discipline.  Likewise, a raw functor
stores functorial action.  A USD functor stores functorial action together
with an adoptive face and a non-adoptive face that are stance-lifts of
the same carrier, together with a boundary witness and a coherence
contract.  Forgetting these fields returns the raw functor.
Keeping them enables safe composition of records without silently turning
records into proof steps.

% -----------------------------------------------------------------------------
\section{Base data}
% -----------------------------------------------------------------------------

This section fixes the data that later sections use.  The goal is to
avoid introducing objects such as $\Rel$, $\BVis$, $\Obs$, or $\NonAdopt$
as undefined prose.

\begin{assumption}[Context category]
Let $\mathcal P$ be a small category.  Its objects are positions,
contexts, observers, stages, regions, or local situations.  A morphism
$\phi:p\to q$ is an admissible transition of context.  In concrete
models $\mathcal P$ may be a poset, a site, or a category of local
charts.
\end{assumption}

\begin{assumption}[Indexed entity category]
Let
\[
  E:\mathcal P^{\op}\to \Cat
\]
be an indexed category.  For each context $p$ the category $E(p)$
contains the entities, expressions, records, or local objects available
at $p$.  For each transition $\phi:p\to q$, the reindexing functor
$E(\phi):E(q)\to E(p)$ gives the local view of a $q$-object at $p$.
\end{assumption}

\begin{construction}[Situated occurrence category]
The category of situated occurrences is the Grothendieck construction
\[
  \Occ(E)\defeq \int_{\mathcal P}E.
\]
An object is a pair $(p,x)$ with $x\in E(p)$.  A morphism
$(p,x)\to(q,y)$ is a pair $(\phi,\alpha)$ where $\phi:p\to q$ in
$\mathcal P$ and
\[
  \alpha:x\to E(\phi)(y)
\]
in $E(p)$.  Composition is the usual composition in the Grothendieck
construction.
\end{construction}

This definition repairs the ambiguity that appears when a set of
positions is used while morphisms $p\to q$ are also required.  From this
point on, context transitions are morphisms in $\mathcal P$.

\begin{assumption}[Relation witness profunctor]
Let $\mathcal C$ be a category.  A relation witness structure on
$\mathcal C$ is a profunctor valued in posets,
\[
  \Rel_{\mathcal C}:\mathcal C^{\op}\times\mathcal C\to \Pos.
\]
An element $r\in \Rel_{\mathcal C}(a,b)$ is a witness that $a$ and $b$
are related in the specified sense.  Functoriality gives transport of
relation witnesses along morphisms of $\mathcal C$.
\end{assumption}

\begin{assumption}[Boundary visibility]
Boundary visibility is a substructure of relation witnesses.  We write
\[
  \BVis_{\mathcal C}\hookrightarrow \Rel_{\mathcal C}
\]
for a subfunctor or indexed subposet.  A boundary witness over
$r\in\Rel_{\mathcal C}(a,b)$ is denoted
\[
  \beta\in \BVis_{\mathcal C}(r).
\]
Thus $\BVis$ is not a metaphor; it is a typed predicate over relation
witnesses.
\end{assumption}

\begin{remark}[Minimality of the base]
Nothing in the formal development requires $\Rel$ to be equality,
isomorphism, equivalence, metric proximity, causal accessibility, or
logical entailment.  Those are model choices.  The USD reference
implementation only assumes relation witnesses and boundary witnesses.
\end{remark}

% -----------------------------------------------------------------------------
\section{The USD record fibration}
% -----------------------------------------------------------------------------

The adoptive and non-adoptive sides are not introduced as two unrelated
categories.  They arise as fibers of a single record fibration.

\begin{definition}[Stance category]
The stance category $\St$ has two distinguished objects
\[
  U,\quad D,
\]
called the adoptive stance and the non-adoptive stance.  We assume an
involutive equivalence
\[
  \tau:\St\to\St,
  \qquad \tau(U)=D,
  \qquad \tau(D)=U,
  \qquad \tau^2\cong 1_{\St}.
\]
The category $\St$ may be taken discrete in the minimal model.  In richer
models it may contain comparison morphisms between stances.
\end{definition}

\begin{definition}[USD record fibration]
A USD record fibration is a functor
\[
  \pi:\RecUSD\to \St.
\]
The fiber over $U$ is the adoptive record category
\[
  \UFib\defeq \pi^{-1}(U),
\]
and the fiber over $D$ is the non-adoptive record category
\[
  \DFib\defeq \pi^{-1}(D).
\]
The notation $\UFib$ and $\DFib$ therefore denotes fibers of one record
fibration, not two arbitrary background categories.
\end{definition}

\begin{definition}[Carrier-realized USD record fibration]
Fix a carrier category $\mathcal C$.  A carrier-realized USD record
fibration over $\mathcal C$ is a USD record fibration together with a
realization functor
\[
  r:\RecUSD\to\mathcal C
\]
and hence a combined projection
\[
  q=(r,\pi):\RecUSD\to\mathcal C\times\St.
\]
We require
\[
  \operatorname{pr}_{\St}\circ q=\pi,
  \qquad
  \operatorname{pr}_{\mathcal C}\circ q=r.
\]
Thus every USD record has both a stance and a realized carrier.  For
$c\in\mathcal C$ and $s\in\St$, the notation
\[
  (\RecUSD)_{(c,s)}\defeq q^{-1}(c,s)
\]
denotes the category of stance-$s$ records realizing the carrier datum
$c$.  The notation suppresses the dependence of $\RecUSD$ on
$\mathcal C$ when no confusion is possible.
\end{definition}

\begin{remark}[Why realization is required]
The projection $\pi:\RecUSD\to\St$ separates adoptive from non-adoptive
records, but by itself it does not say which ordinary carrier object a
record is about.  The realization functor $r:\RecUSD\to\mathcal C$ fills
this gap.  It prevents the pure carrier, the adoptive face, and the
non-adoptive face of a USD functor from becoming three unrelated
functors.
\end{remark}

\begin{definition}[Record constructors]
The record category contains typed constructors
\[
  \Obs,
  \qquad \Adopt,
  \qquad \NonAdopt,
  \qquad \Bnd,
  \qquad \Collapse,
\]
with the following intended types:
\begin{align*}
  \Obs &: \text{raw datum}\to \RecUSD,\\
  \Adopt &: \Rel_{\mathcal C}(a,b)\to \UFib,\\
  \NonAdopt &: \Rel_{\mathcal C}(a,b)\to \DFib,\\
  \Bnd &: \BVis_{\mathcal C}(r)\to \RecUSD,\\
  \Collapse &: \Rel_{\mathcal C}(a,b)\to \RecUSD.
\end{align*}
These constructors are part of the syntax of records.  A model may
interpret them as evidence, annotations, proof-search traces, logs, or
semantic states.  The realization functor $r:\RecUSD\to\mathcal C$
records which carrier datum such a constructed record is about; two
records may have different stances while realizing the same carrier.
\end{definition}

\begin{definition}[Directed stance notation]
Given $x,y$ in a common context, the notation
\[
  x\uconn y
\]
means that a boundary-visible connection record from $x$ to $y$ has been
constructed.  The notation
\[
  x\unadopt y
\]
means that a proposed connection from $x$ to $y$ has been observed but is
recorded on the non-adoptive side.  Neither notation is equality or
ordinary inequality.
\end{definition}

% -----------------------------------------------------------------------------
\section{USD functors as smart functors}
% -----------------------------------------------------------------------------

We now define the main object.  The definition is deliberately close to a
reference implementation: a USD functor is a record with fields.

\begin{definition}[Pure carrier]
Let $\mathcal X$ and $\mathcal C$ be categories.  A pure carrier is an
ordinary functor
\[
  F:\mathcal X\to\mathcal C.
\]
It is the raw functor that remains after forgetting all USD metadata.
\end{definition}


\begin{definition}[Stance-lifts of a carrier]
Let $q=(r,\pi):\RecUSD\to\mathcal C\times\St$ be a carrier-realized USD
record fibration and let $F:\mathcal X\to\mathcal C$ be a pure carrier.
Write
\[
  \underline U,\underline D:\mathcal X\to\St
\]
for the constant functors at the two stances.  A strict adoptive
stance-lift of $F$ is a functor
\[
  U_F:\mathcal X\to\RecUSD
\]
such that
\[
  q\circ U_F=(F,\underline U),
\]
and a strict non-adoptive stance-lift of $F$ is a functor
\[
  D_F:\mathcal X\to\RecUSD
\]
such that
\[
  q\circ D_F=(F,\underline D).
\]
Equivalently, $\pi U_F=\underline U$, $\pi D_F=\underline D$, and
$rU_F=F=rD_F$.  In a lax model, these equalities may be replaced by
specified comparison 2-cells, but the strict form is the reference
implementation used in the main text.
\end{definition}

\begin{definition}[USD functor]
A USD functor from $\mathcal X$ to $\mathcal C$ is a tuple
\[
  \mathbb H=(F,U_F,D_F,\partial_F,\kappa_F)
\]
where:
\begin{enumerate}[label=(\roman*),leftmargin=*]
  \item $F:\mathcal X\to\mathcal C$ is a pure carrier functor.
  \item $U_F:\mathcal X\to\RecUSD$ is a strict adoptive stance-lift of
  $F$, so $q\circ U_F=(F,\underline U)$.
  \item $D_F:\mathcal X\to\RecUSD$ is a strict non-adoptive stance-lift
  of $F$, so $q\circ D_F=(F,\underline D)$.
  \item $\partial_F$ is boundary data over the common carrier $F$,
  relating the two stance-lifts.  A minimal implementation treats
  $\partial_F$ as a family of witnesses
  \[
    \partial_F(x)\in \Bnd_F(U_F(x),D_F(x))
  \]
  natural in $x\in\mathcal X$.
  \item $\kappa_F$ is a coherence contract governing lifts, boundary data,
  morphisms, compositions, erasure, and descent.
\end{enumerate}
The critical point is that $U_F$ and $D_F$ do not merely accompany $F$:
they are records over the same realized carrier $F$.
\end{definition}

\begin{definition}[Coherence contract]
For a USD functor $\mathbb H=(F,U_F,D_F,\partial_F,\kappa_F)$, the
coherence contract is a tuple of obligations
\[
  \kappa_F=(\kappa_F^{\mathrm{lift}},
            \kappa_F^{\partial},
            \kappa_F^{\mathrm{mor}},
            \kappa_F^{\mathrm{comp}},
            \kappa_F^{\mathrm{erase}},
            \kappa_F^{\mathrm{desc}}).
\]
Here:
\begin{enumerate}[label=(\roman*),leftmargin=*]
  \item $\kappa_F^{\mathrm{lift}}$ asserts the stance-lift equations
  $q\circ U_F=(F,\underline U)$ and $q\circ D_F=(F,\underline D)$.
  \item $\kappa_F^{\partial}$ asserts that $\partial_F$ is defined over
  the common realized carrier $F$ and is natural in $\mathcal X$.
  \item $\kappa_F^{\mathrm{mor}}$ specifies what a USD morphism must
  preserve: carrier, U-face, D-face, boundary witness, and contract data.
  \item $\kappa_F^{\mathrm{comp}}$ states the closure obligations for
  $\CComp$, $\DComp$, $\EComp$, $\SComp$, $\TComp$, and $\UComp$.
  \item $\kappa_F^{\mathrm{erase}}$ prevents unlicensed USD records from
  yielding USD-free conclusions after erasure.
  \item $\kappa_F^{\mathrm{desc}}$ records which descent licenses, if any,
  are available for translating records into base assertions.
\end{enumerate}
Thus $\kappa_F$ is not auxiliary prose; it is the contract that makes the
smart functor safe to reference, transform, compose, erase, and descend.
\end{definition}

\begin{definition}[Forgetful functor]
There is a forgetful functor
\[
  |-|:\Smart_{\USD}(\mathcal X,\mathcal C)\to [\mathcal X,\mathcal C]
\]
which sends
\[
  (F,U_F,D_F,\partial_F,\kappa_F)\mapsto F.
\]
This is the categorical version of dereferencing the smart functor to its
pure carrier.
\end{definition}

\begin{remark}[Why this is not merely gluing two categories]
The data $\UFib$ and $\DFib$ are not assumed independently and then
attached to $F$.  They are stance fibers of the single record fibration
$\pi:\RecUSD\to\St$, and the same records are realized in the carrier
category by $r:\RecUSD\to\mathcal C$.  Consequently the adoptive and
non-adoptive faces of a USD functor are stance-lifts over the same pure
carrier:
\[
  q\circ U_F=(F,\underline U),
  \qquad
  q\circ D_F=(F,\underline D).
\]
The boundary witness $\partial_F$ and coherence contract $\kappa_F$
therefore relate two readings of one carrier, not two unrelated functors.
\end{remark}

\begin{definition}[Free and anchored USD functors]
A free USD functor is a USD functor considered before anchoring it to
particular entities.  An anchored USD functor is a free USD functor
equipped with endpoints or a local occurrence context.  We write
\[
  \Anch_{a,b}(\mathbb H)=\mathbb H_{a,b}:a\Rightarrow_{\USD}b
\]
for anchoring $\mathbb H$ between $a$ and $b$ when the required endpoint
and boundary data exist.
\end{definition}

The distinction is essential.  Functions, expressions, programs,
formulas, and hypotheses may be free USD functors before they touch any
particular pair of entities.  Entity-to-entity contact is an anchoring of
an already typed USD functor.

% -----------------------------------------------------------------------------
\section{Morphisms and the category of USD functors}
% -----------------------------------------------------------------------------

\begin{definition}[Morphism of USD functors]
Let
\[
  \mathbb H=(F,U_F,D_F,\partial_F,\kappa_F),\qquad
  \mathbb K=(G,U_G,D_G,\partial_G,\kappa_G)
\]
be USD functors from $\mathcal X$ to $\mathcal C$.  A morphism
$\eta:\mathbb H\to\mathbb K$ consists of:
\begin{enumerate}[label=(\roman*),leftmargin=*]
  \item a natural transformation $\eta^0:F\Rightarrow G$ of pure carriers;
  \item a natural transformation $\eta^U:U_F\Rightarrow U_G$ in
  $[\mathcal X,\RecUSD]$ whose realization is the carrier transformation
  and whose stance component is fixed:
  \[
    q\eta^U=(\eta^0,\mathrm{id}_{\underline U});
  \]
  \item a natural transformation $\eta^D:D_F\Rightarrow D_G$ in
  $[\mathcal X,\RecUSD]$ whose realization is the carrier transformation
  and whose stance component is fixed:
  \[
    q\eta^D=(\eta^0,\mathrm{id}_{\underline D});
  \]
  \item compatibility squares showing that $\eta^U$ and $\eta^D$ preserve
  the boundary data $\partial_F$ and $\partial_G$ over $\eta^0$;
  \item preservation of the coherence contracts, in particular
  $\kappa^{\mathrm{lift}}$, $\kappa^{\partial}$, and
  $\kappa^{\mathrm{erase}}$.
\end{enumerate}
\end{definition}

\begin{proposition}[Category of USD functors]
For fixed $\mathcal X$ and $\mathcal C$, USD functors and their morphisms
form a category, denoted
\[
  \Smart_{\USD}(\mathcal X,\mathcal C),
\]
provided the total record category admits vertical composition of
$q$-compatible natural transformations and the coherence contracts are
closed under identity and composition.
\end{proposition}

\begin{proof}
Identities are inherited componentwise from the functor category and the
record category.  Composition is componentwise and remains compatible
with $q$ because carrier and stance components compose in
$\mathcal C\times\St$.  The boundary compatibility and contract
preservation conditions are stable by the assumed closure of contracts.  The forgetful projection to $[\mathcal X,\mathcal C]$ is
functorial by construction.
\end{proof}

This proposition is intentionally modest.  It is not advertised as a deep
category-theoretic theorem.  It is the sanity check that the reference
implementation has the shape of a category.

% -----------------------------------------------------------------------------
\section{S-molecules and residue-stable sideways composites}
% -----------------------------------------------------------------------------

Before defining the six causeton types, we isolate one record shape that
will be used repeatedly in examples: an \(S\)-molecule.  This section does
not define one of the six causetons and does not state a new proof rule.
It fixes the local data needed to speak about two USD records that remain
distinguishable while sharing a stable boundary-visible discrepancy.  The term is internal
to USD theory; it does not mean a chemical molecule unless a separate
physical realization or descent license is supplied.

\begin{definition}[S-compatibility]
Let \(A\) and \(B\) be USD records or USD functors.  They are
S-compatible when the following data are available:
\begin{enumerate}[label=(\roman*),leftmargin=*]
  \item their realized carriers lie in a common carrier category, or have
  first been brought to a common comparison context;
  \item the relevant relation witnesses and boundary-visibility predicates
  are defined for cross-face comparisons between \(A\) and \(B\);
  \item their coherence contracts permit a shared boundary record and a
  residue record without detaching either U-face or D-face from its
  realized carrier.
\end{enumerate}
S-compatibility is only a typing condition.  It does not identify \(A\)
with \(B\), and it does not imply that their gap is zero.
\end{definition}

\begin{definition}[U/D interlock]
Let \(A\) and \(B\) be S-compatible USD records or USD functors.  A
U/D interlock between \(A\) and \(B\) is boundary-visible data witnessing
that the adoptive face of each side can be read against the non-adoptive
face of the other side.  In notation, this is recorded as
\[
  U_A \uconn D_B,
  \qquad
  U_B \uconn D_A,
\]
together with the relevant boundary witnesses and coherence data.
The symbol \(\uconn\) is not equality.  It records a boundary-visible
connection between stance-lifted records over their realized carriers.
\end{definition}

\begin{definition}[Sideways gap datum]
For S-compatible \(A\) and \(B\), a sideways gap datum consists of a
complete lattice
\[
  \Lambda_{A,B}
\]
of possible gap values, a monotone update map
\[
  \Phi_{A,B}:\Lambda_{A,B}\to\Lambda_{A,B},
\]
and a selected fixed point
\[
  \Delta^*=\Phi_{A,B}(\Delta^*).
\]
The order on \(\Lambda_{A,B}\) is chosen by the model; it may mean
``no more discrepant than'', ``no less resolved than'', or another
typed refinement order.  The selected fixed point is part of the datum
unless it is supplied by an independent fixed-point theorem.  Without
such a point there is no stabilized \(S\)-molecule.
\end{definition}

\begin{definition}[S-molecule]
Let \(A\) and \(B\) be S-compatible USD records or USD functors equipped
with a sideways gap datum.  An \(S\)-molecule between \(A\) and \(B\) is a
sideways-stabilized USD record
\[
  \operatorname{SMol}(A,B;\Delta^*)
\]
consisting of the fixed gap \(\Delta^*\), a U/D interlock between \(A\)
and \(B\), and coherence data asserting that \(\Delta^*\) is preserved by
the U-face, D-face, and boundary witness of the composite.  Material not
adopted by this main observable record is stored in the residue ledger
defined next.

Thus an \(S\)-molecule is not the collapse of \(A\) and \(B\) into one
identity.  It is a USD composite in which two distinguishable sides keep a
shared stable discrepancy and become observable as one composite record.
\end{definition}

\begin{definition}[S-molecular residue ledger]
Let
\[
  \mathbb M=\operatorname{SMol}(A,B;\Delta^*)
\]
be an \(S\)-molecule.  A residue ledger for \(\mathbb M\) is a
self-anchored USD functor
\[
  \mathbb W_{\mathbb M}:\mathbb M\Rightarrow_{\USD}\mathbb M
\]
that records the loss, mismatch, delay, non-adoption, and
boundary-visible waste not adopted by the main observable composite.
The ledger is self-anchored because the residue is about the stabilized
composite itself, not a separate USD-free conclusion.
\end{definition}

\begin{definition}[S-molecular residue accounting]
Let
\[
  \mathsf{Qty}:\{A,B,\mathbb M,\mathbb W_{\mathbb M}\}\to\mathcal R
\]
be a USD quantity valuation into an ordered commutative monoid, defined
for the records under discussion.  A residue ledger for
\(\mathbb M=\operatorname{SMol}(A,B;\Delta^*)\) is conservative when
\[
  \mathsf{Qty}(A)+\mathsf{Qty}(B)
  =
  \mathsf{Qty}(\mathbb M)
  +
  \mathsf{Qty}(\mathbb W_{\mathbb M}).
\]
This equation is a bookkeeping policy inside the USD model.  It does not
turn the unadopted residue into a base-language assertion.
\end{definition}

\begin{remark}[Meaning of molecule]
The word ``molecule'' is used here by analogy with stable composite
entities.  A physical covalent molecule, a polymer, a software package
dependency cluster, or a social institution may be modeled as an
\(S\)-molecule only after specifying an interpretation of \(A\), \(B\),
the gap lattice, the update map, the residue ledger, and the relevant
boundary witnesses.  Without such an interpretation,
\(\operatorname{SMol}(A,B;\Delta^*)\) is only a USD record.
\end{remark}

\begin{remark}[Why residue is part of the molecule]
Without the residue ledger, the stabilized composite may look as if the
unadopted part of the interaction has disappeared.  The ledger prevents
that reading.  Every loss or mismatch must remain accountable either
inside the main observable composite or inside
\(\mathbb W_{\mathbb M}\).  This is consistent with the conservative
design of USD theory: unadopted records are not silently converted into
USD-free conclusions.
\end{remark}

\begin{definition}[UD-spin]
Let
\[
  \mathbb M=\operatorname{SMol}(A,B;\Delta^*)
\]
be an \(S\)-molecule, and let \(\xi\) be an external perturbation, force,
load, attack, or environmental input applied to \(\mathbb M\).  A
UD-spin generated by \(\xi\) is a finite or eventually stationary internal
update sequence
\[
  S_0\to S_1\to\cdots\to S_n
\]
in which the perturbation circulates between the U-face and D-face of
\(\mathbb M\) before being absorbed, recorded in
\(\mathbb W_{\mathbb M}\), or ejected through boundary-visible residue.

UD-spin is a USD bookkeeping operation.  It is not physical spin unless a
separate physical model identifies it with a physical spin-like quantity.
\end{definition}

\begin{definition}[Divergence, ejection, and spin termination]
Let \(\omega\) denote a UD-spin speed for an \(S\)-molecule \(\mathbb M\).
Let
\[
  \mathsf{Div}_{\mathbb M}(\omega)
\]
be the amount of perturbation that \(\mathbb M\) can internally dissipate
or diverge at spin speed \(\omega\), and let
\[
  \mathsf{Eject}_{\mathbb M}(\omega)
\]
be the amount that can be ejected through boundary-visible residue at
that speed.  If \(\lambda_\xi\) denotes the incoming perturbation rate,
then the basic stability condition is
\[
  \lambda_\xi
  \leq
  \mathsf{Div}_{\mathbb M}(\omega)
  +
  \mathsf{Eject}_{\mathbb M}(\omega).
\]
When this condition holds and the induced UD-spin terminates, the
\(S\)-molecule returns to a state that is observationally equivalent, at
the chosen boundary resolution, to the original composite record.
\end{definition}

\begin{definition}[Critical spin speed]
The critical spin speed of an \(S\)-molecule \(\mathbb M\) is the boundary
spin speed at which the incoming perturbation rate begins to exceed the
sum of the divergence capacity and ejection capacity:
\[
  \lambda_\xi
  >
  \mathsf{Div}_{\mathbb M}(\omega)
  +
  \mathsf{Eject}_{\mathbb M}(\omega).
\]
Beyond this boundary, residue accumulates faster than it can be processed.
If \(R_t\) denotes accumulated residue, one may record the update
schematically as
\[
  R_{t+1}
  =
  R_t
  +
  \lambda_\xi
  -
  \bigl(
    \mathsf{Div}_{\mathbb M}(\omega)
    +
    \mathsf{Eject}_{\mathbb M}(\omega)
  \bigr).
\]
When the residue exceeds the tolerance allowed by the coherence contract,
the \(S\)-molecule may break, dissociate, reconfigure, or cease to be
observable as one composite record.
\end{definition}

\begin{definition}[Observed UD-spin speed]
The internal UD-spin speed of an \(S\)-molecule and the externally
observed UD-spin speed need not coincide.  We write
\[
  \omega^{\mathrm{int}}_{\mathbb M}
\]
for the internal spin speed and
\[
  \omega^{\Obs}_{\mathbb M}
\]
for the boundary-visible observed spin speed.  If
\(\alpha_{\mathbb M}\) denotes observational resolution or accuracy, then
one may model the observed speed as
\[
  \omega^{\Obs}_{\mathbb M}
  =
  \alpha_{\mathbb M}\cdot \omega^{\mathrm{eff}}_{\mathbb M}.
\]
For a long serial chain of \(S\)-molecules, observational resolution may
decrease as the chain length increases.  In such models the observed
UD-spin speed is allowed to remain constant or decrease, even if local
internal spin speeds do not decrease.
\end{definition}

\begin{definition}[Non-observable pending wait]
During the formation or perturbation of an \(S\)-molecule, internal
non-equivalent exchanges need not be externally observable.  Let
\(\sigma_A\) be the time at which an internal accumulation begins in
\(A\), let \(\theta_A\) be the time at which that accumulation becomes
boundary-visible, and let \(\tau_{A\to B}\) be the time at which the
corresponding exchange is implemented toward \(B\).  The non-observable
pending wait from \(A\) to \(B\) is
\[
  T^{\mathrm{NP}}_{A\to B}
  =
  \tau_{A\to B}-\theta_A.
\]
Here NP means ``non-observable pending'' and is unrelated to the
complexity class NP.
\end{definition}

\begin{definition}[S-polymer]
An \(S\)-polymer is a finite or indexed serial joining of \(S\)-molecules
\[
  \operatorname{SMol}_1
  \mathbin{\otimes_S}
  \operatorname{SMol}_2
  \mathbin{\otimes_S}
  \cdots
  \mathbin{\otimes_S}
  \operatorname{SMol}_n
\]
which is externally observable as one higher-order USD entity at a chosen
boundary resolution.  Each adjacent joining must supply its own
S-compatibility data, sideways gap datum, interlock, and residue ledger.
An \(S\)-polymer may model a physical material, a software dependency
cluster, an institution, or another composite system, provided that a
suitable realization of the USD records is supplied.
\end{definition}

\begin{remark}[S-molecules and physical examples]
A covalent bond may be used as an intuitive model of an \(S\)-molecule:
two atoms remain distinguishable while forming a stable shared state
through electron-density-mediated interaction.  This comparison is only a
model reading.  USD theory does not redefine covalent bonding, molecular
physics, software dependency, or any other base-domain concept.  To obtain
a USD-free physical conclusion, the relevant descent license must still be
supplied in the base theory.
\end{remark}

% -----------------------------------------------------------------------------
\section{Causetons: six composition--decomposition pairs}
% -----------------------------------------------------------------------------

The six basic USD operations are now treated as causetons.  A causeton is
not merely a composition step.  It is a typed record of a composition
operation together with the decomposition phenomenon emitted by that
operation, the self-anchored residue ledger that stores what is not
adopted by the main output, and the coherence contract that prevents the
record from becoming an unlicensed proof rule.

\begin{definition}[Causeton]
For
\[
  \star\in\{C,D,E,S,T,U\},
\]
a \(\star\)-causeton is a typed USD causal record
\[
  \mathfrak c_\star
  =
  (\Comp_\star,\Dec_\star,\mathbb W_\star,\kappa_\star).
\]
Here \(\Comp_\star\) is the \(\star\)-typed composition operation,
\(\Dec_\star\) is the decomposition phenomenon emitted by that operation,
\(\mathbb W_\star\) is a self-anchored residue USD functor recording the
unadopted byproduct, and \(\kappa_\star\) is the coherence contract for the
causeton.
\end{definition}

\begin{remark}[Causetons are USD records]
A causeton is not asserted to be a physical particle, a new causal law, or
a base-language explanation.  It is a USD record unit.  A physical,
historical, legal, or psychological reading requires the relevant carrier,
relation witnesses, boundary witnesses, and descent licenses.
\end{remark}

\begin{assumption}[Common closure invariant for causetons]
Every causeton below is required to preserve carrier-realized
stance-lifts.  If the main output has carrier
\[
  F':\mathcal X'\to\mathcal C,
\]
then its output faces must satisfy
\[
  q\circ U_{F'}=(F',\underline U),
  \qquad
  q\circ D_{F'}=(F',\underline D).
\]
This is the operational content of $\kappa^{\mathrm{comp}}$: composition
may change the carrier, boundary, or indexing shape, and decomposition may
emit residue, but neither side may detach U-face or D-face records from
their carrier realization.
\end{assumption}

\begin{definition}[C-causeton: compression and interface decomposition]
The C-causeton
\[
  \mathfrak c_C
  =
  (\Comp_C,\Dec_C,\mathbb W_C,\kappa_C)
\]
has the following data.

\textbf{Composition operation.}  The input consists of two free USD
functors
\[
  \mathbb H:\mathcal X\to\mathcal C,
  \qquad
  \mathbb K:\mathcal Y\to\mathcal C,
\]
together with an interface span
\[
  \mathcal X\xleftarrow{i}\mathcal I\xrightarrow{j}\mathcal Y
\]
and interface compatibility data between $i^*\mathbb H$ and
$j^*\mathbb K$.  When the pushout
\[
  \mathcal X\sqcup_{\mathcal I}\mathcal Y
\]
exists in $\Cat$, and the carrier and record data glue along the
interface in a way compatible with $q$, the composition operation is
\[
  \Comp_C(\mathbb H,\mathbb K)
  =
  \CComp_{\mathcal I}(\mathbb H,\mathbb K):
  \mathcal X\sqcup_{\mathcal I}\mathcal Y\to\mathcal C.
\]

\textbf{Decomposition phenomenon.}  Compression emits the decomposition
of the interface.  The old boundary data do not disappear; the interface
discrepancy becomes separable as residue.  A minimal boundary policy is
\[
  \partial_{\CComp(\mathbb H,\mathbb K)}
  =
  \partial_{\mathbb H}\oplus
  \partial_{\mathbb K}\oplus
  \Delta_{\mathcal I}(\mathbb H,\mathbb K),
\]
where $\Delta_{\mathcal I}$ is supplied by the model.

\textbf{Residue ledger.}  The self-anchored ledger
\[
  \mathbb W_C:
  \CComp_{\mathcal I}(\mathbb H,\mathbb K)
  \Rightarrow_{\USD}
  \CComp_{\mathcal I}(\mathbb H,\mathbb K)
\]
records the projection loss, interface discrepancy, and non-adopted
boundary material emitted by the compression.
\end{definition}

\begin{definition}[D-causeton: restriction and background decomposition]
The D-causeton
\[
  \mathfrak c_D
  =
  (\Comp_D,\Dec_D,\mathbb W_D,\kappa_D)
\]
has the following data.

\textbf{Composition operation.}  The input is a context morphism
\[
  \rho:\mathcal L\to\mathcal G
\]
and a USD functor over the global context $\mathcal G$.  The composition
operation is reindexing or pullback:
\[
  \Comp_D(\rho,\mathbb H)
  =
  \DComp_{\rho}(\mathbb H)
  =
  \rho^*\mathbb H:
  \mathcal L\to\mathcal C.
\]
If $\mathbb H$ has carrier $F$, then the output has carrier $\rho^*F$,
and its faces are the reindexed stance-lifts satisfying
\[
  q(\rho^*U_F)=(\rho^*F,\underline U),
  \qquad
  q(\rho^*D_F)=(\rho^*F,\underline D).
\]

\textbf{Decomposition phenomenon.}  Downward restriction decomposes the
global record into the local reading and the background material that the
local context no longer displays.  The local output is not a proof that the
global claim holds locally; it is the typed reading of the global record
through the chosen index.

\textbf{Residue ledger.}  The self-anchored ledger
\[
  \mathbb W_D:
  \rho^*\mathbb H\Rightarrow_{\USD}\rho^*\mathbb H
\]
records the forgotten global constraints, boundary conditions, and
nonlocal residue left outside the local restriction.
\end{definition}

\begin{definition}[E-causeton: cut extension and path decomposition]
The E-causeton
\[
  \mathfrak c_E
  =
  (\Comp_E,\Dec_E,\mathbb W_E,\kappa_E)
\]
has the following data.

\textbf{Composition operation.}  The input is an anchored USD functor
\[
  \mathbb H_{a,b}:a\Rightarrow_{\USD}b
\]
and a proposed cut vertex $m$.  When there is factorization data in the
relevant arrow category, together with a comparison 2-cell
\[
  \varepsilon_m:
  \mathbb H_{a,m}\circ \mathbb H_{m,b}
  \Rightarrow
  \mathbb H_{a,b}
\]
or the reverse comparison, the composition operation is
\[
  \Comp_E(\mathbb H_{a,b};m)
  =
  \EComp_m(\mathbb H_{a,b})
  =
  (\mathbb H_{a,m},\mathbb H_{m,b};\varepsilon_m).
\]

\textbf{Decomposition phenomenon.}  The cut vertex decomposes a direct
anchored record into mediated partial records.  The decomposition is a
subdivision or refinement, not a proof that the original record is
strictly equal to the subdivided composite.

\textbf{Residue ledger.}  The self-anchored ledger
\[
  \mathbb W_E:
  \EComp_m(\mathbb H_{a,b})
  \Rightarrow_{\USD}
  \EComp_m(\mathbb H_{a,b})
\]
records the explanatory excess of the cut vertex, the direction of
$\varepsilon_m$, and the material not adopted by either partial path.
\end{definition}

\begin{definition}[S-causeton: stabilization and gap decomposition]
The S-causeton
\[
  \mathfrak c_S
  =
  (\Comp_S,\Dec_S,\mathbb W_S,\kappa_S)
\]
has the following data.

\textbf{Composition operation.}  The input is a pair of S-compatible USD
functors $\mathbb H,\mathbb K$, together with a complete gap lattice
\[
  \Lambda_{\mathbb H,\mathbb K}
\]
and a monotone gap update map
\[
  \Phi_{\mathbb H,\mathbb K}:
  \Lambda_{\mathbb H,\mathbb K}\to
  \Lambda_{\mathbb H,\mathbb K}.
\]
After a fixed point
\[
  \Delta^*=\Phi_{\mathbb H,\mathbb K}(\Delta^*)
\]
is selected, the composition operation is
\[
  \Comp_S(\mathbb H,\mathbb K;\Delta^*)
  =
  \SComp(\mathbb H,\mathbb K;\Delta^*)
  \defeq
  \operatorname{SMol}(\mathbb H,\mathbb K;\Delta^*).
\]
Existence follows, for example, from the Knaster-Tarski fixed point
theorem when the stated monotonicity and completeness conditions hold.

\textbf{Decomposition phenomenon.}  Sideways stabilization decomposes the
interaction into the shared observable molecule and the gap that remains
nonzero.  The completion condition is not $\Delta^*=0$ but shared
stability.

\textbf{Residue ledger.}  The self-anchored ledger is the
S-molecular ledger already introduced:
\[
  \mathbb W_S
  =
  \mathbb W_{\operatorname{SMol}(\mathbb H,\mathbb K;\Delta^*)}:
  \operatorname{SMol}(\mathbb H,\mathbb K;\Delta^*)
  \Rightarrow_{\USD}
  \operatorname{SMol}(\mathbb H,\mathbb K;\Delta^*).
\]
It records the non-adopted discrepancy, fixed sideways gap, and
boundary-visible waste emitted by the stabilization.
\end{definition}

\begin{definition}[Sequence category]
Let $\Seq(\Smart_{\USD})$ be the category of finite composable sequences
of USD functors.  An object is a finite chain
\[
  \mathcal H=(\mathbb H_1,\ldots,\mathbb H_n).
\]
\end{definition}

\begin{definition}[U/D flip]
Assume that the stance involution $\tau:\St\to\St$ lifts to a
carrier-preserving record involution
\[
  \bar\tau:\RecUSD\to\RecUSD,
  \qquad
  r\bar\tau=r,
  \qquad
  \pi\bar\tau=\tau\pi.
\]
The induced operation
\[
  \Flip_{U,D}:\Smart_{\USD}\to\Smart_{\USD}
\]
sends
\[
  (F,U_F,D_F,\partial_F,\kappa_F)
  \mapsto
  (F,\bar\tau D_F,\bar\tau U_F,\partial_F^\tau,\kappa_F^\tau).
\]
The pure carrier is not altered.  The lifted faces satisfy
\[
  q(\bar\tau D_F)=(F,\underline U),
  \qquad
  q(\bar\tau U_F)=(F,\underline D).
\]
\end{definition}

\begin{definition}[T-causeton: turnover and directional decomposition]
The T-causeton
\[
  \mathfrak c_T
  =
  (\Comp_T,\Dec_T,\mathbb W_T,\kappa_T)
\]
has the following data.

\textbf{Composition operation.}  The input is a finite USD functor chain
\[
  \mathcal H=(\mathbb H_1,\ldots,\mathbb H_n).
\]
The composition operation is temporal turnover:
\[
  \Comp_T(\mathcal H)
  =
  \TComp(\mathcal H)
  \defeq
  \Rev\bigl(\Flip_{U,D}(\mathbb H_1),\ldots,
             \Flip_{U,D}(\mathbb H_n)\bigr),
\]
so that
\[
  \TComp(\mathcal H)
  =
  (\mathbb H_n^\tau,\ldots,\mathbb H_1^\tau).
\]

\textbf{Decomposition phenomenon.}  Turnover decomposes the original
reading order, causal direction, responsibility direction, and adopted
stance assignment into a turned record chain.  A chain that was read
forward as an identification proposal may be re-read backward as a
non-adoption record, but this is a transformation of records, not a proof
of ordinary inequality.

\textbf{Residue ledger.}  The self-anchored ledger
\[
  \mathbb W_T:
  \TComp(\mathcal H)\Rightarrow_{\USD}\TComp(\mathcal H)
\]
records the directional residue emitted by reversal and U/D turnover.  If
no information is lost, $\TComp$ is involutive up to record equivalence:
\[
  \TComp^2\cong 1.
\]
If compression, stabilization, or descent has discarded information, the
weaker USD-level comparison
\[
  \TComp^2(\mathcal H)\uconn \mathcal H
\]
may be all that remains.
\end{definition}

\begin{definition}[U-causeton: upward image]
D-composition has two common adjoints when the relevant Kan extensions
exist.  The U-causeton
\[
  \mathfrak c_U
  =
  (\Comp_U,\Dec_U,\mathbb W_U,\kappa_U)
\]
has the following data.

\textbf{Composition operation.}  Let $\rho:\mathcal L\to\mathcal G$ be a
context morphism.  When the left Kan extension exists, define the
existential upward operation by
\[
  \Comp_U^{\exists}
  =
  \UComp^{\exists}_{\rho}
  \defeq
  \Sigma_\rho
  \defeq
  \Lan_{\rho}:
  \Smart_{\USD}(\mathcal L,\mathcal C)\to
  \Smart_{\USD}(\mathcal G,\mathcal C),
\]
with
\[
  \Sigma_\rho\dashv \rho^*.
\]
When the right Kan extension exists, define the universal upward operation
by
\[
  \Comp_U^{\forall}
  =
  \UComp^{\forall}_{\rho}
  \defeq
  \Pi_\rho
  \defeq
  \Ran_{\rho}:
  \Smart_{\USD}(\mathcal L,\mathcal C)\to
  \Smart_{\USD}(\mathcal G,\mathcal C),
\]
with
\[
  \rho^*\dashv \Pi_\rho.
\]
These Kan extensions are admitted as U-causetons only when they lift from
the pure carrier category to $\Smart_{\USD}$.  For example, if
$\Sigma_\rho\mathbb H$ has carrier $\Lan_\rho F$, then its faces must be
records over that carrier:
\[
  q\circ U_{\Lan_\rho F}=(\Lan_\rho F,\underline U),
  \qquad
  q\circ D_{\Lan_\rho F}=(\Lan_\rho F,\underline D),
\]
and similarly for $\Pi_\rho$ and $\Ran_\rho F$.

\textbf{Decomposition phenomenon.}  Upward image decomposes a local record
into the global candidate and the local specificity not adopted by that
global reading.  The existence of the ordinary Kan extension of $F$ is not
enough; the USD record and contract data must lift with it.

\textbf{Residue ledger.}  The self-anchored ledger
\[
  \mathbb W_U:
  \UComp_{\rho}(\mathbb H)\Rightarrow_{\USD}\UComp_{\rho}(\mathbb H)
\]
records the boundary conditions, local assumptions, and non-generalized
material left behind by the upward image.  If the lifted Kan extension
does not exist, the corresponding U-causeton is undefined in the smart
category even when the carrier Kan extension exists.
\end{definition}

\subsection{Causeton coverage and saturation}

The six causetons are not phases of an algorithm.  A model may instantiate
them in any order, may instantiate the same causeton type more than once,
and may feed the records produced by one causeton into any other causeton
whenever the input and existence conditions match.  What can be required
without ambiguity is not a global application order, but a coverage
certificate.

\begin{definition}[Causeton agenda]
For a USD test model, a causeton agenda is a six-tuple
\[
  \mathcal A
  =
  (\mathcal A_C,\mathcal A_D,\mathcal A_E,
   \mathcal A_S,\mathcal A_T,\mathcal A_U),
\]
where $\mathcal A_\star$ is a finite or indexed family of typed candidate
instances of the corresponding causeton.  The agenda is
six-causeton-covered when no component is empty and every listed candidate
supplies the input data,
existence conditions, decomposition phenomenon, residue ledger, and
coherence contract required by its causeton type.
\end{definition}

\begin{definition}[Saturation closure]
Let $R_0$ be a family of seed USD records and let $\mathcal A$ be a
six-causeton-covered causeton agenda.  Write
\[
  \Gamma_{\mathcal A}(R)
  =
  R\cup C_{\mathcal A}(R)\cup D_{\mathcal A}(R)
   \cup E_{\mathcal A}(R)\cup S_{\mathcal A}(R)
   \cup T_{\mathcal A}(R)\cup U_{\mathcal A}(R)
\]
for the family obtained by adding all well-typed outputs and residue
ledgers of the candidate instances whose inputs are present in $R$.  A
saturated expansion is a fixed point
\[
  R^*=\Gamma_{\mathcal A}(R^*)
\]
generated from $R_0$; in a finite agenda this is the least closure reached
by repeatedly adding newly well-typed outputs.  The expansion is accepted
only when every generated record preserves the common stance-lift
invariant through $q$.
\end{definition}

\begin{definition}[Causeton coverage certificate]
A coverage certificate for a USD test model consists of a
six-causeton-covered causeton agenda $\mathcal A$, a saturated expansion
$R^*$, and a check that the realized causeton types are exactly
\[
  \{\mathfrak c_C,\mathfrak c_D,\mathfrak c_E,
    \mathfrak c_S,\mathfrak c_T,\mathfrak c_U\}.
\]
Thus missing use of any one causeton type is a failure of the test model,
while multiple instances of a causeton type are permitted whenever they are
typed.
\end{definition}

Diagrams in examples may still display dependency arrows between records.
Such arrows describe which outputs and residue ledgers are used as later
inputs; they are not a global execution order for the six causetons.

% -----------------------------------------------------------------------------
\section{Sensitons: art-anchored perceptual connectors}
% -----------------------------------------------------------------------------

Causetons record typed composition--decomposition pairs.  Sensitons record
how selected causeton precedence patterns are anchored inside an artistic
observation context and become perceptible as controlled abstraction.  A
sensiton is therefore not a violation pattern.  It is a connector from a
causeton pair to an observer's cognitive S-polymer under the boundary of a
work, medium, performance, recording, exhibition, or other art-anchored
frame.

\begin{definition}[Art anchoring]
An art anchoring is a context
\[
  \Art
\]
whose objects are works, media, performances, recordings, scripts, scores,
images, scenes, audiences, exhibition spaces, or other records belonging
to an artistic observation frame.  The anchoring supplies a boundary
contract asserting that records inside \(\Art\) are not automatically
licensed as USD-free real-world commands, empirical claims, or moral
judgments.
\end{definition}

\begin{remark}[Why art anchoring is needed]
The precedence patterns below can be unsafe when they are read as
unlicensed claims about the base world.  Inside \(\Art\), the same
patterns may become aesthetic operations: symbol formation, plot
compression, reversal, surprise, estrangement, catharsis, or controlled
recomposition.  Art anchoring is the contract that keeps the emitted
fission perceptible without immediately turning it into an external
descent.
\end{remark}

\begin{definition}[Sensiton]
Let \(O\) be an observer with cognitive S-polymer \(\mathbb P_O\).  A
sensiton is an art-anchored USD perceptual record
\[
  \mathfrak s
  =
  (\mathfrak c_X\prec\mathfrak c_Y,
   \Art,
   O,
   \mathbb P_O,
   \mathbb W_{\mathfrak s},
   \kappa_{\mathfrak s}),
\]
where \(\mathfrak c_X\prec\mathfrak c_Y\) means that the \(X\)-causeton
is perceptually staged before the \(Y\)-causeton inside \(\Art\).  The
ledger \(\mathbb W_{\mathfrak s}\) records the perceptual residue emitted
by that staging, and \(\kappa_{\mathfrak s}\) keeps the record inside the
art anchoring unless an explicit descent license is supplied.
\end{definition}

\begin{definition}[Magic-Sensiton]
The Magic-Sensiton is the art-anchored precedence pattern
\[
  \MagicSens
  \defeq
  \Sens(U\prec D).
\]
It stages a U-causeton before the corresponding D-causeton.  In an
unanchored setting this may look like an upward image before local
restriction, context, or verification.  Inside \(\Art\), the same pattern
is read as magic, symbol, mythic ascent, or imaginative world-formation:
a local or weakly anchored record is first lifted into a larger possible
world and only afterward lowered into scenes, bodies, materials, or
particular perceptions.
\end{definition}

\begin{definition}[Plot-Sensiton]
The Plot-Sensiton is the art-anchored precedence pattern
\[
  \PlotSens
  \defeq
  \Sens(C\prec E).
\]
It stages a C-causeton before the corresponding E-causeton.  In an
unanchored setting this may compress anchored records before the relevant
cut vertices, tests, or mediation points are exposed.  Inside \(\Art\),
the same pattern is read as plot: fragments are first compressed into a
line, myth, scene, motif, or narrative pressure, and only afterward
subdivided by turns, cuts, revelations, delays, and reversals.
\end{definition}

\begin{definition}[Hall-Sensiton]
The Hall-Sensiton is the art-anchored precedence pattern
\[
  \HallSens
  \defeq
  \Sens(T\prec S).
\]
It stages a T-causeton before the corresponding S-causeton.  In an
unanchored setting this may reverse a system before its stabilizing
S-molecular reading has been supplied.  Inside \(\Art\), the same pattern
is read as hall, stage, theater, or performance space: time order,
viewpoint, role, and U/D stance are turned before the work restabilizes
them as an observable scene.  The observer may experience surprise,
estrangement, dread, recognition, or catharsis, but these effects remain
anchored to the work unless a descent license carries them outside
\(\Art\).
\end{definition}

\begin{remark}[Order of the three basic sensitons]
The order
\[
  \MagicSens,\qquad
  \PlotSens,\qquad
  \HallSens
\]
follows the increasing externalization of abstraction.  Magic-Sensiton
forms a possible world, Plot-Sensiton compresses fragments into a line of
readability, and Hall-Sensiton stages the line before an observer through
turned time, role, and stance.
\end{remark}

\begin{remark}[Sensitons and USD fission]
Sensitons are designed to emit controlled USD fission in the observer's
cognitive S-polymer.  The fission is not treated as a defect merely
because it occurs.  It becomes aesthetically useful when the art anchoring
allows the observer to hold the fragments, residues, and reversed readings
long enough for recomposition inside the work.
\end{remark}

% -----------------------------------------------------------------------------
\section{USD fusion and USD fission}
% -----------------------------------------------------------------------------

Causetons separate the typed operation from what an observer sees.  The
operation side is composition and decomposition.  The observational side
is USD fusion and USD fission.  Fusion and fission are treated here as two
provisional USD phenomena; later refinements may classify them into
C/D/E/S/T/U-specific fusion and fission types.

\begin{definition}[USD fusion]
Let \(\mathfrak c_\star\) be a causeton with main output \(\mathbb M\),
residue ledger \(\mathbb W_\star\), and observer or boundary resolution
\(O\).  USD fusion is the observer-relative event
\[
  \USDFusion_O(\mathfrak c_\star,\mathbb M)
\]
in which \(\mathbb M\) is readable as one boundary-visible composite
record.  Fusion is not identity collapse: the residue ledger remains
attached, and the U-face and D-face remain stance-lifted over their
realized carrier.
\end{definition}

\begin{definition}[USD fission]
Let \(\mathfrak c_\star\) be a causeton with decomposition phenomenon
\(\Dec_\star\), residue ledger \(\mathbb W_\star\), and observer or entity
polymer \(\mathbb P_O\).  USD fission is the observer-relative event
\[
  \USDFission_O(\mathfrak c_\star,\mathbb P_O)
\]
in which the emitted decomposition becomes readable inside \(\mathbb P_O\)
as multiple fragments or separated partial records.  It may also appear as
contradictions, unintelligible residue, uncanny residue, or other
boundary-visible non-unified material.  Fission is a USD observation, not
by itself a USD-free assertion that the source theory is false.
\end{definition}

\begin{definition}[Controlled and critical fission]
USD fission is controlled for an observer \(O\) when the emitted fragments
can be reassembled, interpreted, or used under the observer's coherence
contract without losing the relevant boundary witnesses.  It is critical
when the emitted residue exceeds the observer's available divergence,
ejection, or recomposition capacity and the record can no longer be read as
one stable composite at the chosen resolution.
\end{definition}

\begin{remark}[No intrinsic moral valence]
USD fusion and USD fission are observational phases, not moral judgments.
A fission may be destructive, controlled, explanatory, analytic, or
creative depending on the surrounding carrier model and descent licenses.
Likewise, a fusion may be stabilizing, misleading, compressive, or
over-adoptive depending on the residue ledger it emits.
\end{remark}

\begin{remark}[Composition and decomposition versus fusion and fission]
Composition and decomposition belong to the causeton record.  Fusion and
fission are how the output and emitted residue are observed.  Thus the
reference slogan is:
\[
  \text{composition and fusion, decomposition and fission}.
\]
This slogan does not identify USD fusion with physical nuclear fusion, nor
USD fission with physical nuclear fission.  Those readings require
separate base-domain realization.
\end{remark}

\section{A syntactic reference implementation}
% -----------------------------------------------------------------------------

This section records the minimum syntax needed for a conservative USD
calculus.  It is intentionally austere.

\begin{definition}[Base language]
Let $\mathcal L_0$ be a base formal language with judgments
\[
  \Gamma\vdash_0 A.
\]
The base language may be first-order logic, type theory, an internal
language of a category, or another fixed formal system.
\end{definition}

\begin{definition}[USD extension]
The USD calculus $\USDCalc$ extends $\mathcal L_0$ with record terms:
\[
  \mathrm{conn}(x,y;r,\beta),
  \qquad
  \mathrm{nonadopt}(x,y;r,\beta),
  \qquad
  \mathrm{hand}(F,U,D,\partial,\kappa),
\]
and with record-forming operations
\[
  \CComp,
  \DComp,
  \EComp,
  \SComp,
  \TComp,
  \UComp.
\]
The USD extension does not add a base rule allowing a record term alone to
prove a USD-free formula.
\end{definition}

\begin{definition}[Descent license]
A descent license is an explicit rule instance
\[
  \IdLic(R,A)
\]
that permits a USD record $R$ to be translated into a base assertion $A$
in $\mathcal L_0$.  Examples include a definition, an isomorphism, a
proved equivalence, a semantic interpretation, or a verified coercion.
\end{definition}

\begin{definition}[Erasure]
The erasure translation
\[
  |-|:\USDCalc\to\mathcal L_0
\]
acts as identity on base terms and formulas and removes pure USD records
unless a descent license supplies a base formula.
\end{definition}

\begin{theorem}[Conservative design theorem]
Assume that every USD inference rule is record-forming, record-transforming,
or explicitly licensed by a descent rule.  If
\[
  \Gamma\vdash_{\USDCalc} A
\]
and $A$ is USD-free, then the derivation of $A$ uses only base rules and
licensed descents.  Consequently
\[
  |\Gamma|\vdash_0 A
\]
provided each used descent license is valid in the base system.
\end{theorem}

\begin{proof}
By induction on the given derivation.  Base rules translate to base rules.
Record-forming and record-transforming USD rules have USD-record
conclusions and therefore cannot be the final step of a USD-free formula.
If a USD record is used to produce a USD-free conclusion, the rule must be
a licensed descent by hypothesis; replacing that step with its base
license yields a derivation in $\mathcal L_0$.  The induction removes all
unlicensed USD bookkeeping from the derivation of $A$.
\end{proof}

This theorem is deliberately conditional.  It states the safety contract
of the reference implementation.  A stronger metatheorem for a particular
formal calculus would require a fully specified proof system, but this
condition is already enough to prevent USD records from acting as hidden
proof rules.

% -----------------------------------------------------------------------------
\section{Worked toy model}
% -----------------------------------------------------------------------------

\begin{example}[Poset record model]
Let $\mathcal C$ be a poset regarded as a category.  Let
\[
  \Rel_{\mathcal C}(a,b)
\]
be the truth value of $a\leq b$.  Let $\BVis$ select those comparable
pairs for which a chosen boundary label is present.  Let $\RecUSD$ be the
product category
\[
  \mathcal C\times\St
\]
with projection to $\St$.  The realization projection is
\[
  r=\operatorname{pr}_{\mathcal C}:\mathcal C\times\St\to\mathcal C,
\]
so the combined projection is
\[
  q=(r,\pi):\mathcal C\times\St\to\mathcal C\times\St,
\]
the identity in this minimal model.  Then
$\UFib\cong\mathcal C\cong\DFib$, but the two copies are distinguished as
fibers.  A USD functor is a functor $F$ together with two stance-lifted
copies $U_F,D_F$ over the same $F$ and a boundary label.  In this model
D-composition is ordinary reindexing of monotone maps, C-composition is
pushout of indexing posets when it exists, and T-composition swaps the
two copies and reverses a finite chain while preserving realization.
\end{example}

\begin{example}[Complete lattice gap model]
Let the gap lattice be a complete lattice $\Lambda$.  If a sideways
interaction gives a monotone map $\Phi:\Lambda\to\Lambda$, then the set
of fixed points is nonempty by the Knaster-Tarski theorem.  Choosing the
least fixed point gives a canonical minimal shared discrepancy.  Choosing
the greatest fixed point gives a maximal conservative discrepancy.  The
choice is part of the model policy.
\end{example}

% -----------------------------------------------------------------------------
\section{Relativistic test model}
% -----------------------------------------------------------------------------

The relativistic material is not an application claiming new physics.  It
is a test model for the USD vocabulary because relativity naturally
involves observers, frames, local restrictions, and time-oriented records.
The historical sources are Einstein's 1905 paper on special relativity
and his 1916 paper on general relativity \cite{Einstein1905,Einstein1916}.

\begin{definition}[Lorentzian test base]
Let $(M,g)$ be a Lorentzian manifold.  Let $\mathcal O(M,g)$ be a chosen
category of local observers or frames.  A morphism in $\mathcal O(M,g)$
may be a change of frame, inclusion of a local chart, or another chosen
admissible comparison of observations.
\end{definition}

\begin{example}[D-composition as local restriction]
A global geometric record, such as a metric-dependent record on $(M,g)$,
may be lowered along a local frame map
\[
  \rho:O\to M
\]
by pullback:
\[
  \DComp_{\rho}(\mathbb H)=\rho^*\mathbb H.
\]
This is the familiar mathematical shape of restricting global data to a
local observer or frame.  The USD interpretation only adds adoptive and
non-adoptive record faces and boundary witnesses.
\end{example}

\begin{example}[T-composition as record turnover]
A time-oriented chain of observation records may be transformed by
$\TComp$ into the reverse-order chain with U and D faces exchanged.  This
is not asserted to be physical time reversal symmetry.  It is a record
operation that leaves the entities and geometric carrier untouched while
turning the USD bookkeeping upside down.
\end{example}

\begin{remark}[No hidden physics]
If a model uses Lorentz transformations, connections, curvature, or field
equations, those structures must be explicitly included in the base
category or in the relation witness profunctor.  USD functors do not
supply them automatically.
\end{remark}

\subsection{Newton-Einstein test}

The Newton-Einstein test uses E-composition as its main operation.  The
test does not claim that Einstein proved Newtonian absolute time false by
one direct replacement.  It records a two-stage theoretical operation:
Newton's single global time parameter is first decomposed inside special
relativity and then extended through local inertial structure into
general relativity.  The relevant historical base texts are Newton's
\emph{Principia}, Einstein's 1905 paper on special relativity, and
Einstein's 1916 paper on general relativity
\cite{Newton1687,Einstein1905,Einstein1916}.

Let
\[
  \mathcal N_{\mathrm{mech}},
  \qquad
  \mathcal S_{\mathrm{rel}},
  \qquad
  \mathcal G_{\mathrm{rel}}
\]
denote schematic categories of Newtonian mechanics, special relativity,
and general relativity.  A Newtonian time record may be written
\[
  \mathbb N_T:a\Rightarrow_{\USD} b,
\]
where an event configuration is assigned to a single frame-independent
temporal order.  In the Newtonian record, simultaneity and duration are
treated as globally available, and the time coordinate is not produced by
a synchronization protocol internal to the theory.

The coarse E-composition flow is
\[
  \mathbb N_T
  \xrightarrow{\EComp_{\mathcal S_{\mathrm{rel}}}}
  \mathbb E_{\mathrm{SR}}
  \xrightarrow{\EComp_{\mathcal G_{\mathrm{rel}}}}
  \mathbb E_{\mathrm{GR}}.
\]
Thus special relativity is not skipped.  It is the first cut vertex
through which the Newtonian record is factored before the general
relativistic record is formed.

\textbf{First E-instance: Newton to special relativity.}  The cut vertex
$m_{\mathrm{SR}}$ is itself resolved into smaller cuts:
\[
\begin{aligned}
  \mathbb N_T
  &\xrightarrow{\EComp_{\mathrm{clock}}}
  \text{local clock readings}\\
  &\xrightarrow{\EComp_{\mathrm{sync}}}
  \text{light-signal synchronization}\\
  &\xrightarrow{\EComp_{\mathrm{inertial}}}
  \text{inertial-frame coordinate time}\\
  &\xrightarrow{\EComp_{\mathrm{Lorentz}}}
  \text{Lorentz-compatible transformations}\\
  &\xrightarrow{\EComp_{\mathrm{interval}}}
  \mathbb E_{\mathrm{SR}}.
\end{aligned}
\]
The central cut is synchronization.  What was transparent in the
Newtonian record becomes an explicit operation:
\[
  \text{remote simultaneity}
  \quad\leadsto\quad
  \text{clock comparison by signal protocol}.
\]
The carrier is no longer a single absolute time assignment.  It is a
special-relativistic carrier in which coordinate time is indexed by
inertial frame and constrained by the constancy of light speed.  The
adoptive face contains
\[
\begin{aligned}
  U_{\mathrm{SR}}
  =
  \{&\text{local clocks, synchronization, inertial frames,}\\
    &\text{Lorentz transformations, invariant interval}\}.
\end{aligned}
\]
The non-adoptive residue contains
\[
\begin{aligned}
  D_{\mathrm{Newton}}
  =
  \{&\text{universal simultaneity, frame-independent absolute time,}\\
    &\text{Galilean velocity addition, instantaneous global temporal order}\}.
\end{aligned}
\]
This is not erasure of Newtonian mechanics.  It is a factorization of the
role played by absolute time into observable clock readings, synchronization,
frame choice, transformation law, and invariant spacetime interval.

\textbf{Second E-instance: special to general relativity.}  The next cut
vertex is not another synchronization convention.  It is the passage from
global inertial frames to local inertial structure:
\[
\begin{aligned}
  \mathbb E_{\mathrm{SR}}
  &\xrightarrow{\EComp_{\mathrm{local}}}
  \text{local inertial frames}\\
  &\xrightarrow{\EComp_{\mathrm{equiv}}}
  \text{equivalence principle}\\
  &\xrightarrow{\EComp_{g}}
  \text{metric field }g\\
  &\xrightarrow{\EComp_{\mathrm{curv}}}
  \text{curvature}\\
  &\xrightarrow{\EComp_{\mathrm{EFE}}}
  \mathbb E_{\mathrm{GR}}.
\end{aligned}
\]
Special relativity remains as a local record: in sufficiently small
regions and in the absence of curvature effects at the chosen resolution,
the general relativistic record restricts to a special-relativistic one.
The general relativistic extension changes the carrier from flat
Minkowski spacetime to metric spacetime.  Time is no longer merely
frame-indexed coordinate time on a flat background; it is part of a
metric field whose geometry is coupled to gravitation.

The complete theoretical surgery can therefore be summarized as
\[
\begin{aligned}
  \text{absolute Newtonian time}
  &\xrightarrow{\EComp_{\mathcal S_{\mathrm{rel}}}}
  \text{clock/synchronization/frame/interval structure}\\
  &\xrightarrow{\EComp_{\mathcal G_{\mathrm{rel}}}}
  \text{local inertial/metric/curvature structure}.
\end{aligned}
\]
The USD reading is that Einstein did not merely negate the Newtonian
record.  He inserted the missing operational vertices, preserved the
usable low-speed and weak-field residue through appropriate limiting
policies, and moved the privileged invariant from absolute time to
spacetime structure.  A USD-free physical claim about the Newtonian limit,
Lorentz transformations, or Einstein field equations still requires the
ordinary base theory; the USD test only records the typed decomposition
flow.

% -----------------------------------------------------------------------------
\section{Interdisciplinary test model}
% -----------------------------------------------------------------------------

The preceding relativistic section tests USD vocabulary inside a
mathematical physics setting.  The same discipline can be used for
interdisciplinary examples, where the carrier data are not all drawn from
one field.  The purpose of such a model is not to turn a social or
economic interpretation into a proof of intention.  It is to keep the
interfaces visible when social records and economic records are read
together.

\begin{definition}[Interdisciplinary test base]
Let $\mathcal S$ be a category of sociological records and let
$\mathcal E$ be a category of economic records.  An interdisciplinary
carrier may be taken to be the product category
\[
  \mathcal C_{\mathcal S,\mathcal E}
  \defeq
  \mathcal S\times\mathcal E,
\]
or, more generally, a category equipped with projection functors
\[
  p_{\mathcal S}:\mathcal C_{\mathcal S,\mathcal E}\to\mathcal S,
  \qquad
  p_{\mathcal E}:\mathcal C_{\mathcal S,\mathcal E}\to\mathcal E.
\]
A USD functor into $\mathcal C_{\mathcal S,\mathcal E}$ records a
joint social-economic reading.  Its relation witnesses may compare
social data, economic data, or boundary-visible couplings between them.
\end{definition}

\begin{remark}[No hidden interdisciplinary descent]
An interdisciplinary USD record is not, by itself, an empirical
explanation.  A descent license must specify the accepted historical
sources, social model, economic model, and rule that permits a USD record
to be translated into a USD-free claim.  Without such a license, the
record remains a structured hypothesis or interpretation.
\end{remark}

\subsection{Satoshi test}

The pseudonymous creator of Bitcoin gives a useful interdisciplinary
test case.  The Bitcoin white paper presents the system as peer-to-peer
electronic cash without reliance on a trusted third party
\cite{Nakamoto2008}.  Late in 2010, Satoshi Nakamoto publicly warned that
Bitcoin was still a small beta community and that sudden WikiLeaks-driven
attention could bring destructive heat \cite{Nakamoto2010WikiLeaks}.
The following day, Satoshi also described denial-of-service resistance as
unfinished work \cite{Nakamoto2010DoS}.  Separately, claims that Satoshi
or a dominant early miner controlled about one million bitcoins depend on
chain-analysis heuristics; the evidence is suggestive but not as robust
as popular retellings often imply \cite{BitMEX2018}.

The source data are kept in two different categories.  The sociological
side records legitimacy, founder authority, regulatory visibility,
community dependence, and governance centrality.  The economic side
records early coin concentration, liquidity, expected sell-pressure
signals, scarcity belief, and market confidence.  The test is not
licensed to derive an economic conclusion directly from a sociological
one, or a sociological intention directly from an economic trace.

Let $\mathcal A_{\mathcal S}$ be a small category of sociological founder
actions: revealing an identity, exercising founder authority, transferring
control, remaining pseudonymous, or withdrawing from public control.  Let
$\mathcal A_{\mathcal E}$ be a small category of economic founder actions:
moving early coins, not moving early coins, distributing control,
signalling sale, or remaining economically inactive.  These categories are
not identified.  They touch only through a boundary-visible interface.

\textbf{D0-instance.}  Begin with broad sociological and economic USD records
$\mathbb H_{\mathcal S}^{0}$ and $\mathbb H_{\mathcal E}^{0}$.  Let
\[
  \rho_{\mathcal S}^{0}:
  \mathcal A_{\mathcal S}\to\mathcal G_{\mathcal S},
  \qquad
  \rho_{\mathcal E}^{0}:
  \mathcal A_{\mathcal E}\to\mathcal G_{\mathcal E}
\]
be context morphisms selecting only founder-action data from broader
social and economic contexts.  The first downward restriction is
\[
  \mathbb H_{\mathcal S}^{f}
  \defeq
  \DComp_{\rho_{\mathcal S}^{0}}(\mathbb H_{\mathcal S}^{0}),
  \qquad
  \mathbb H_{\mathcal E}^{f}
  \defeq
  \DComp_{\rho_{\mathcal E}^{0}}(\mathbb H_{\mathcal E}^{0}).
\]
This step does not select withdrawal as the answer.  It only removes
material outside the founder-action comparison class.

\textbf{C-instance.}  The social and economic records are read together only
inside the interdisciplinary carrier $\mathcal C_{\mathcal S,\mathcal E}$.
Before applying compression, assume the model supplies typed carrier lifts
\[
  \widehat{\mathbb H}_{\mathcal S}^{f}:
  \mathcal A_{\mathcal S}\to \mathcal C_{\mathcal S,\mathcal E},
  \qquad
  \widehat{\mathbb H}_{\mathcal E}^{f}:
  \mathcal A_{\mathcal E}\to \mathcal C_{\mathcal S,\mathcal E}
\]
whose projections recover the restricted social and economic records:
\[
  p_{\mathcal S}\circ\widehat{\mathbb H}_{\mathcal S}^{f}
  =
  \mathbb H_{\mathcal S}^{f},
  \qquad
  p_{\mathcal E}\circ\widehat{\mathbb H}_{\mathcal E}^{f}
  =
  \mathbb H_{\mathcal E}^{f}.
\]
The unused projections are bookkeeping fields supplied by the model, not
licenses for cross-disciplinary descent.
Let
\[
  \Theta_{\mathrm{Sat}}:
  \mathcal A_{\mathcal S}^{\op}\times\mathcal A_{\mathcal E}
  \to \Pos
\]
be an interface profunctor.  An element
\[
  \theta\in\Theta_{\mathrm{Sat}}(s,e)
\]
witnesses that a social founder action $s$ and an economic founder action
$e$ are two boundary-visible readings of the same founder episode, such as
public identity, authority, early coin movement, non-movement, market
signal, or community dependence.  Let
\[
  \mathcal I_{\mathrm{Sat}}
  \defeq
  \int_{\mathcal A_{\mathcal S}^{\op}\times\mathcal A_{\mathcal E}}
  \Theta_{\mathrm{Sat}}.
\]
Compression along this interface gives the joint record
\[
  \mathbb J_{\mathrm{Sat}}
  \defeq
  \CComp_{\mathcal I_{\mathrm{Sat}}}
  (\widehat{\mathbb H}_{\mathcal S}^{f},
   \widehat{\mathbb H}_{\mathcal E}^{f}).
\]
The compression does not say that social authority is an economic price,
or that economic non-movement is a social intention.  It records only the
typed interface through which both projections become visible in the same
interdisciplinary carrier.

\textbf{D1-instance.}  Now restrict the compressed record to the local Satoshi
context.  Let
\[
  \rho_{\mathrm{Sat}}:
  \mathcal L_{\mathrm{Sat}}
  \to
  \mathcal A_{\mathcal S}\sqcup_{\mathcal I_{\mathrm{Sat}}}
  \mathcal A_{\mathcal E}
\]
select the local span generated by pseudonymity, public withdrawal,
non-exercise of founder authority, early coin non-movement, and the
visible decentralization signal.  Define
\[
  \mathbb L_{\mathrm{Sat}}
  \defeq
  \DComp_{\rho_{\mathrm{Sat}}}(\mathbb J_{\mathrm{Sat}}).
\]
This second downward restriction is local reading, not historical
descent.  It does not infer the inner intention, the identity, or the key
status of Satoshi.

\textbf{E-instance.}  Anchor the local record between an initial founder-action
state $a_0$ and a decentralization-signal state $a_1$:
\[
  \mathbb L_{\mathrm{Sat},a_0,a_1}:a_0\Rightarrow_{\USD}a_1.
\]
Insert the following cut vertices by iterated E-composition:
\[
\begin{aligned}
  m_{\mathrm{id}} &= \text{public identity visibility},\\
  m_{\mathrm{auth}} &= \text{founder authority},\\
  m_{\mathrm{coin}} &= \text{early coin movement or non-movement},\\
  m_{\mathrm{sell}} &= \text{large sell-pressure signal}.
\end{aligned}
\]
Thus
\[
  \mathbb E_{\mathrm{Sat}}
  \defeq
  \EComp_{m_{\mathrm{sell}}}
  \EComp_{m_{\mathrm{coin}}}
  \EComp_{m_{\mathrm{auth}}}
  \EComp_{m_{\mathrm{id}}}
  (\mathbb L_{\mathrm{Sat},a_0,a_1}).
\]
The first two cuts are sociological: they concern visibility, authority,
central target formation, and community dependence.  The last two cuts are
economic: they concern coin movement and sell-pressure signals.  The
factorization is deliberately conservative; it prevents the record from
jumping directly from withdrawal to decentralization or from coin
non-movement to personal intention.

\textbf{S-instance.}  The sideways stabilization is the core of the Satoshi
test.  Let $\Lambda_{\mathrm{Sat}}$ be a complete lattice of discrepancies
between decentralization-supporting outputs and non-adopted founder
residue, and let
\[
  \Phi_{\mathrm{Sat}}:
  \Lambda_{\mathrm{Sat}}\to\Lambda_{\mathrm{Sat}}
\]
be the monotone update map supplied by the model.  Choose a fixed point
\[
  \Delta_{\mathrm{Sat}}^*
  =
  \Phi_{\mathrm{Sat}}(\Delta_{\mathrm{Sat}}^*).
\]
The adoptive face is a product reading
\[
  U_{\mathrm{decent}}
  =
  (U_{\mathcal S},U_{\mathcal E}),
\]
where
\[
\begin{aligned}
  U_{\mathcal S}
  &=
  \{\text{reduced founder authority, reduced central target,
    reduced community dependence}\},\\
  U_{\mathcal E}
  &=
  \{\text{reduced large-sell signal, reduced direct founder-liquidity
    pressure}\}.
\end{aligned}
\]
The non-adoptive face is likewise a product reading
\[
  D_{\mathrm{fame/wealth}}
  =
  (D_{\mathcal S},D_{\mathcal E}),
\]
where
\[
\begin{aligned}
  D_{\mathcal S}
  &=
  \{\text{public fame, founder control, proved inner intention}\},\\
  D_{\mathcal E}
  &=
  \{\text{direct realization of early coin wealth, exact status of
    early coins}\}.
\end{aligned}
\]
Let $A_{\mathrm{decent}}$ be the decentralization-support record whose
adoptive face is $U_{\mathrm{decent}}$, and let
$B_{\mathrm{residue}}$ be the founder-residue record whose non-adoptive
face is $D_{\mathrm{fame/wealth}}$.
The stabilized record is
\[
  \operatorname{SMol}_{\mathrm{Sat}}
  \defeq
  \SComp(A_{\mathrm{decent}},B_{\mathrm{residue}};
   \Delta_{\mathrm{Sat}}^*)
  =
  \operatorname{SMol}
  (A_{\mathrm{decent}},B_{\mathrm{residue}};
   \Delta_{\mathrm{Sat}}^*).
\]
At the visible boundary,
\[
  U_{\mathrm{decent}}
  \uconn
  D_{\mathrm{fame/wealth}}.
\]
This says that decentralized adoption can be modeled as stable only while
public fame, founder control, direct realization of early coin wealth,
proved intention, and exact early-coin status remain non-adopted residue.
It does not say that fame or wealth were morally rejected, and it does
not prove who Satoshi was.

A conservative best-effort policy can be expressed by assigning each
compatible local action $a\in\mathcal L_{\mathrm{Sat}}$ a social cost and
an economic cost,
\[
  C_{\mathcal S}(a),
  \qquad
  C_{\mathcal E}(a),
\]
and selecting an action whose worst visible cost is minimal:
\[
  a^*
  \in
  \operatorname*{arg\,min}_{a\in\mathcal L_{\mathrm{Sat}}}
  \max\{C_{\mathcal S}(a),C_{\mathcal E}(a)\}.
\]
Under one plausible policy, withdrawal together with non-movement of early
coins is a candidate for $a^*$.  The social reading is reduced founder
centrality and reduced central target formation.  The economic reading is
the absence of a large direct sell-pressure signal.  The loss is not
erased: abandoned fame, unresolved identity, founder control, direct
realization of early wealth, and exact early-coin status remain attached
as self-anchored residue.

\textbf{T-instance.}  Let $\mathcal H_{\mathrm{Sat}}$ be the finite subdivided
record chain underlying $\mathbb E_{\mathrm{Sat}}$ and
$\operatorname{SMol}_{\mathrm{Sat}}$.  Temporal turnover gives
\[
  \mathbb T_{\mathrm{Sat}}
  \defeq
  \TComp(\mathcal H_{\mathrm{Sat}}).
\]
Forward, the record reads:
\[
  \text{founder withdrawal and early-coin non-movement}
  \longrightarrow
  \text{decentralization-supporting signal}.
\]
Turned over, it reads:
\[
\begin{aligned}
  &\text{decentralized Bitcoin record}\\
  &\quad\longrightarrow
  \text{founder fame, control, wealth, identity, and intention as residue}.
\end{aligned}
\]
This reversal blocks a common overreading.  The present decentralized
record may be read as carrying founder-residue, but it does not reveal
the founder's private reason for withdrawing or not moving coins.

\textbf{U-instance.}  When the relevant lifted Kan extension exists,
image the turned local record upward along $\rho_{\mathrm{Sat}}$:
\[
  \mathbb U_{\mathrm{Sat}}
  \defeq
  \UComp^{\exists}_{\rho_{\mathrm{Sat}}}(\mathbb T_{\mathrm{Sat}})
  \leadsto
  \mathrm{FounderExitAsDecentralizationStabilizer}.
\]
The upward image is a general schema, not a biographical claim.  It says
that, in projects where a founder has strong initial authority and a large
visible economic position, public withdrawal, non-exercise of authority,
and non-movement of large early holdings can stabilize a decentralization
record, provided the fame, control, wealth, identity, and intention that
are not adopted by the main output remain recorded as residue.

\begin{remark}[No proof of Satoshi's intention]
The Satoshi test licenses at most a structural conclusion: withdrawal and
coin non-movement can be modeled as a best-effort stable record for
Bitcoin's social and economic decentralization.  It does not license the
USD-free assertion that Satoshi chose withdrawal for that reason.  The
inner intention, the exact identity, and the exact status of the early
coins remain non-adopted residue unless supplied by independent historical
evidence and an explicit descent license.
\end{remark}

\subsection{Oppenheimer test}

The Oppenheimer test is an interdisciplinary stress test for ethical
translation.  Its source data use only two categories: a category of
theoretical nuclear physics and a category of public policy.  The
philosophical forum and the trial of Socrates are not assumed as source
categories.  They appear only after the test supplies a six-causeton
coverage certificate and saturates the generated USD record family.  No
global application order of the causetons is part of the claim.

Historically, the Manhattan Project made nuclear fission and explosive
chain reactions into an organized weapon program.  The Smyth Report
describes the scientific and administrative development of atomic energy
for military purposes \cite{Smyth1945}.  The Franck Report records that
some project scientists understood the use of the first atomic bombs as a
social and political problem, not merely a technical one
\cite{Franck1945}.  Plato's \emph{Apology} records the trial in which
Socrates' public practice of examination was answered by charges of
corrupting the youth and impiety \cite{PlatoApology}.  The USD claim
below is not that these events are historically the same.  It is that
the Oppenheimer record can be saturated, without dropping the decisive
gap, into a record family whose upward image realizes the forum structure
made visible by the trial of Socrates.

Let
\[
  \mathcal T_{\mathrm{nuc}}
\]
be a category of theoretical nuclear-physics records.  Its objects may be
fission models, chain-reaction calculations, criticality estimates,
yield predictions, uncertainty bounds, or experimental confirmations.
Its morphisms are refinements, derivations, approximations, empirical
constraints, or transitions from speculative theory to reliable
prediction.  Let
\[
  \mathcal P_{\mathrm{pub}}
\]
be a category of public-policy records.  Its objects may be nonresearch,
research, secret development, demonstration, military use, nonuse,
international control, or postwar governance.  Its morphisms are funding,
authorization, classification, command, deployment, diplomatic proposal,
or legal constraint.

The two categories touch through a feasibility profunctor
\[
  \Theta_{\mathrm{Op}}:
  \mathcal T_{\mathrm{nuc}}^{\op}\times
  \mathcal P_{\mathrm{pub}}
  \to \Pos.
\]
An element
\[
  \theta\in\Theta_{\mathrm{Op}}(t,p)
\]
is a witness that a theoretical record $t$ makes a policy record $p$
available, urgent, dangerous, or governable.  Such a witness may record
technical feasibility, predictability, destructive scale, uncertainty,
time pressure, enemy-access risk, secrecy, irreversibility, or
accountability.  The interface category is the Grothendieck construction
\[
  \mathcal I_{\Theta}
  \defeq
  \int_{\mathcal T_{\mathrm{nuc}}^{\op}\times
        \mathcal P_{\mathrm{pub}}}
  \Theta_{\mathrm{Op}},
\]
whose objects are triples $(t,p,\theta)$.  This interface is the exact
place where theoretical success becomes a public-policy option.

\textbf{C-instance.}  Let $\mathbb H_T$ be a USD record over
$\mathcal T_{\mathrm{nuc}}$ and let $\mathbb H_P$ be a USD record over
$\mathcal P_{\mathrm{pub}}$.  To type the compression, assume an
interdisciplinary carrier category $\mathcal C_{\mathrm{Op}}$ and typed
carrier lifts
\[
  \widehat{\mathbb H}_T:
  \mathcal T_{\mathrm{nuc}}\to\mathcal C_{\mathrm{Op}},
  \qquad
  \widehat{\mathbb H}_P:
  \mathcal P_{\mathrm{pub}}\to\mathcal C_{\mathrm{Op}}.
\]
Let
\[
  i:\mathcal I_{\Theta}\to\mathcal T_{\mathrm{nuc}},
  \qquad
  j:\mathcal I_{\Theta}\to\mathcal P_{\mathrm{pub}}
\]
be the interface projections, with the first projection read through the
variance of $\Theta_{\mathrm{Op}}$.  Compression along the interface gives
an Oppenheimer record
\[
  \mathbb O
  \defeq
  \CComp_{\mathcal I_{\Theta}}
  (\widehat{\mathbb H}_T,\widehat{\mathbb H}_P).
\]
Its compressed global context is
\[
  \mathcal G_{\mathrm{Op}}
  \defeq
  \mathcal T_{\mathrm{nuc}}
  \sqcup_{\mathcal I_{\Theta}}
  \mathcal P_{\mathrm{pub}}.
\]
The compression does not identify theory with policy.  It only records
that the two sides now share a boundary witness: the same nuclear-physics
result that counts as explanatory success in $\mathcal T_{\mathrm{nuc}}$
also counts as a possible state action in $\mathcal P_{\mathrm{pub}}$.

\textbf{D-instance.}  Choose a local theory-practice context
\[
  \rho_{\mathrm{Op}}:
  \mathcal L_{\mathrm{Op}}\to\mathcal G_{\mathrm{Op}}
\]
that retains the possible/ought collision and forgets the details of
isotopes, laboratories, committees, and military offices.  Downward
restriction gives
\[
  \mathbb L_{\mathrm{Op}}
  \defeq
  \DComp_{\rho_{\mathrm{Op}}}(\mathbb O).
\]
The local carrier retains the typed collision
\[
  \text{possible}
  \;\Longrightarrow\;
  \text{must decide whether it ought to be done}.
\]
This is sharper than the slogan that can does not imply ought.  The
lowered record says that can creates a decision burden.  A theory may
remain epistemic in its own category, but after contact with public
policy it generates an actionable option for which someone becomes
answerable.

\textbf{E-instance.}  To avoid a leap from theory to guilt, anchor
$\mathbb L_{\mathrm{Op}}$ between the local records
$\mathrm{possible}$ and $\mathrm{responsible}$:
\[
  \mathbb L_{\mathrm{Op},\mathrm{possible},\mathrm{responsible}}:
  \mathrm{possible}\Rightarrow_{\USD}\mathrm{responsible}.
\]
Insert an
internal cut vertex
\[
  m_{\mathrm{act}}
  =
  \text{actionability}.
\]
The Oppenheimer record is subdivided as
\[
  \mathbb E_{\mathrm{Op}}
  \defeq
  \EComp_{m_{\mathrm{act}}}
  (\mathbb L_{\mathrm{Op},\mathrm{possible},\mathrm{responsible}})
  =
  (\mathbb O_{\mathrm{possible},m_{\mathrm{act}}},
   \mathbb O_{m_{\mathrm{act}},\mathrm{responsible}};
   \varepsilon_{m_{\mathrm{act}}}).
\]
The first half records the passage from theoretical possibility to
available state action.  The second half records the passage from
available state action to responsibility for irreversible consequences.
Thus the ethical problem is not imported from outside physics or outside
policy.  It is produced at the cut where a prediction becomes an
implementable public act.

\textbf{S-instance.}  Let $\Lambda_{\mathrm{Op}}$ be a complete lattice
of discrepancies between theoretical success and ethically non-neutral
public action.  Let
\[
  \Phi_{\mathrm{Op}}:
  \Lambda_{\mathrm{Op}}\to\Lambda_{\mathrm{Op}}
\]
be the monotone update map supplied by the model, and choose a fixed
point
\[
  \Delta_{\mathrm{Op}}^*
  =
  \Phi_{\mathrm{Op}}(\Delta_{\mathrm{Op}}^*).
\]
Let $A_{\mathrm{theory}}$ be the theoretical-success record and
$B_{\mathrm{practice}}$ the public-action residue record.  The stable
Oppenheimer molecule is
\[
  \operatorname{SMol}_{\mathrm{Op}}
  \defeq
  \SComp(A_{\mathrm{theory}},B_{\mathrm{practice}};
   \Delta_{\mathrm{Op}}^*)
  =
  \operatorname{SMol}
  (A_{\mathrm{theory}},B_{\mathrm{practice}};
   \Delta_{\mathrm{Op}}^*).
\]
Its adoptive theoretical face is
\[
  U_{\mathrm{theory}}
  =
  \{\text{truth, prediction, possibility, technical control}\},
\]
while its non-adoptive public-policy face is
\[
  D_{\mathrm{practice}}
  =
  \{\text{harm, authority, secrecy, irreversibility,
    accountability}\}.
\]
The interlock is boundary-visible:
\[
  U_{\mathrm{theory}}
  \uconn
  D_{\mathrm{practice}}.
\]
The fixed gap is not zero.  It is the stable discrepancy
\[
\begin{aligned}
  \Delta_{\mathrm{Op}}^*
  &=
  \text{the stable gap between theoretical success}\\
  &\quad\text{and ethically non-neutral public action}.
\end{aligned}
\]
The residue contains both sides of the policy dilemma: use may create
direct and irreversible destruction, while nonuse, delay, demonstration,
or secrecy may leave war, arms-race, or governance risks unresolved.
The molecule therefore preserves the dilemma instead of resolving it.

\textbf{T-instance.}  Let $\mathcal H_{\mathrm{Op}}$ be a finite
composable USD chain containing the subdivided local record
$\mathbb E_{\mathrm{Op}}$ and the stabilized molecule
$\operatorname{SMol}_{\mathrm{Op}}$ as typed factors.  Temporal turnover
gives
\[
  \mathbb T_{\mathrm{Op}}
  \defeq
  \TComp(\mathcal H_{\mathrm{Op}}).
\]
This reverses the reading of the record chain.
Forward, the Oppenheimer molecule reads:
\[
  \text{theory discovers possibility}
  \longrightarrow
  \text{policy must decide action}.
\]
Turned over, it reads:
\[
  \text{public consequence}
  \longrightarrow
  \text{theoretical practice is judged by what it enabled}.
\]
The turned record has the same carrier gap, but its vocabulary can now be
read as a forum problem.  In the forum normal form, theoretical discovery
is replaced by contradiction discovery, and public-policy action is
replaced by public authority over the speaker:
\[
  \text{contradiction is found}
  \;\Longrightarrow\;
  \text{the forum must decide whether to rewind or expel}.
\]
This is the no-rewind forum dilemma.  A truth-seeking argument treats
contradiction as a repair point.  An authority-protecting forum may treat
the same contradiction as disqualification.  The boundary witness is the
same shape as in the Oppenheimer molecule: a successful epistemic act
creates an irreversible practical burden.

\textbf{U-instance.}  Let
\[
  \rho_{\mathrm{Forum}}:
  \mathcal L_{\mathrm{no\text{-}rewind}}
  \to
  \mathcal G_{\mathrm{Forum}}
\]
be the context morphism from the local no-rewind forum pattern to a
global forum schema.  When the relevant lifted Kan extension exists, the
upward image of the turned record is the Socratic trial schema:
\[
  \mathbb U_{\mathrm{Op}}
  \defeq
  \UComp^{\exists}_{\rho_{\mathrm{Forum}}}(\mathbb T_{\mathrm{Op}})
  \leadsto
  \mathrm{Trial}_{\mathrm{Socrates}}.
\]
In this image, the theoretical side becomes Socratic examination:
questions expose contradictions and force interlocutors back to their
premises.  The public-policy side becomes the authority of the Athenian
forum: the city can treat the exposed instability not as an invitation to
rewind the argument, but as corruption, impiety, or civic danger.  The
mapping is:
\[
\begin{array}{ccl}
  \text{nuclear possibility} &\mapsto& \text{exposed contradiction},\\
  \text{state action} &\mapsto& \text{public judgment},\\
  \text{irreversible use} &\mapsto& \text{irreversible punishment},\\
  \text{policy responsibility} &\mapsto& \text{forum authority},\\
  \text{unresolved residue} &\mapsto& \text{loss of rewind rights}.
\end{array}
\]
The coverage certificate for the Oppenheimer test is the
six-causeton-covered agenda
\[
\begin{aligned}
  \mathcal A_C
  &=
  \{\CComp_{\mathcal I_{\Theta}}
    (\widehat{\mathbb H}_T,\widehat{\mathbb H}_P)\},\\
  \mathcal A_D
  &=
  \{\DComp_{\rho_{\mathrm{Op}}}(\mathbb O)\},\\
  \mathcal A_E
  &=
  \{\EComp_{m_{\mathrm{act}}}
    (\mathbb L_{\mathrm{Op},\mathrm{possible},\mathrm{responsible}})\},\\
  \mathcal A_S
  &=
  \{\SComp(A_{\mathrm{theory}},B_{\mathrm{practice}};
    \Delta_{\mathrm{Op}}^*)\},\\
  \mathcal A_T
  &=
  \{\TComp(\mathcal H_{\mathrm{Op}})\},\\
  \mathcal A_U
  &=
  \{\UComp^{\exists}_{\rho_{\mathrm{Forum}}}
    (\mathbb T_{\mathrm{Op}})\}.
\end{aligned}
\]
The agenda may be saturated in any order, and the displayed dependencies
only say which generated records are available as inputs to other typed
causeton instances.  A stronger historical or moral conclusion is not
produced by the saturation itself.
The transformation is therefore not a metaphorical jump from physics to
Athens.  It passes through the common S-molecular form:
\[
  \text{epistemic success}
  \uconn
  \text{practical authority over irreversible consequence}.
\]

\begin{remark}[No proof of historical identity]
The Oppenheimer test licenses a structural comparison, not a USD-free
historical identity.  It does not say that Oppenheimer was Socrates, that
the Manhattan Project was a trial, or that Athenian law and nuclear
policy were the same institution.  It says that both records can realize
the same no-rewind pattern: a practice whose internal norm is correction
by truth is recoded by public authority as a source of irreversible
judgment.  Any stronger historical or moral conclusion requires an
explicit descent license.
\end{remark}

% -----------------------------------------------------------------------------
\section{Reviewer-facing checklist}
% -----------------------------------------------------------------------------

For reference, the implementation satisfies the following design checks.

\begin{enumerate}[leftmargin=*]
  \item The context parameter is a category $\mathcal P$, not a bare set.
  \item Situated occurrences are defined by a Grothendieck construction.
  \item $\Rel$ is a poset-valued profunctor or indexed predicate.
  \item $\BVis$ is a typed substructure of $\Rel$.
  \item $\UFib$ and $\DFib$ are fibers of one fibration
  $\RecUSD\to\St$, and every record is also realized by
  $r:\RecUSD\to\mathcal C$.
  \item The combined projection $q=(r,\pi):\RecUSD\to\mathcal C\times\St$
  makes U- and D-faces stance-lifts over one carrier.
  \item A USD functor is a tuple
  $(F,U_F,D_F,\partial_F,\kappa_F)$ where $U_F$ and $D_F$ satisfy
  $q\circ U_F=(F,\underline U)$ and $q\circ D_F=(F,\underline D)$.
  \item The coherence contract $\kappa_F$ is an explicit bundle of lift,
  boundary, morphism, composition, erasure, and descent obligations.
  \item C/D/E/S/T/U operations have stated input types, outputs, and
  existence conditions.
  \item S-composition depends on a complete lattice and a monotone update
  map, so fixed-point existence is not merely asserted.
  \item T-composition is sequence reversal plus a carrier-preserving U/D
  involution and does not touch entities or carrier realization.
  \item USD-free conclusions require base reasoning or an explicit
  descent license.
\end{enumerate}

% -----------------------------------------------------------------------------
\section{Open problems}
% -----------------------------------------------------------------------------

The present draft is a reference implementation rather than a final
metatheory.  The main open problems are:
\begin{enumerate}[leftmargin=*]
  \item Refine preferred subclasses of coherence contracts $\kappa_F$ and
  prove closure under all six causeton composition--decomposition types in
  nontrivial models.
  \item Compare the carrier-realized USD record fibration with displayed categories,
  double categories, fibrations of proofs, and decorated functors.
  \item Give a fully formal sequent calculus or type theory for
  $\USDCalc$ and prove a cut-elimination or model-theoretic conservativity
  theorem.
  \item Identify useful model policies for selecting fixed points in
  S-composition.
  \item Develop lax stance-lift variants where $q\circ U_F$ and $q\circ D_F$ compare to
  $(F,\underline U)$ and $(F,\underline D)$ rather than strictly equal them.
  \item Develop nontrivial test models beyond the poset and relativistic
  sketches.
\end{enumerate}

% -----------------------------------------------------------------------------
\section{Conclusion}
% -----------------------------------------------------------------------------

The central correction in this draft is that USD Functor Theory should be
read as a specification of smart functors, not as a metaphorical use of
category theory.  A USD functor is an ordinary functor with additional
record fields: adoptive face, non-adoptive face, boundary witness, and
coherence contract.  The U and D regions are fibers of one stance-indexed
record fibration, and the faces are stance-lifts realized over the same
carrier by
\[
  q=(r,\pi):\RecUSD\to\mathcal C\times\St.
\]
The six causetons are typed composition--decomposition pairs backed by
ordinary categorical constructions or explicit existence assumptions, and
each causeton is defined only when it preserves the stance-lift invariant
through $q$.  The resulting system can guide interpretation and proof
search while remaining conservative over the base language for USD-free
conclusions.

\begin{thebibliography}{99}

\bibitem{MacLane1971}
Saunders Mac Lane.
\newblock \emph{Categories for the Working Mathematician}.
\newblock Springer, 1971.

\bibitem{Borceux1994}
Francis Borceux.
\newblock \emph{Handbook of Categorical Algebra, Volumes 1--3}.
\newblock Cambridge University Press, 1994.

\bibitem{Johnstone2002}
Peter T. Johnstone.
\newblock \emph{Sketches of an Elephant: A Topos Theory Compendium}.
\newblock Oxford University Press, 2002.

\bibitem{GrothendieckSGA1}
Alexander Grothendieck.
\newblock Rev\^etements \`etales et groupe fondamental.
\newblock In \emph{SGA 1}, Lecture Notes in Mathematics 224, Springer, 1971.

\bibitem{Tarski1955}
Alfred Tarski.
\newblock A lattice-theoretical fixpoint theorem and its applications.
\newblock \emph{Pacific Journal of Mathematics}, 5(2):285--309, 1955.

\bibitem{Newton1687}
Isaac Newton.
\newblock \emph{Philosophi\ae{} Naturalis Principia Mathematica}.
\newblock 1687.

\bibitem{Einstein1905}
Albert Einstein.
\newblock Zur Elektrodynamik bewegter K\"orper.
\newblock \emph{Annalen der Physik}, 322(10):891--921, 1905.

\bibitem{Einstein1916}
Albert Einstein.
\newblock Die Grundlage der allgemeinen Relativit\"atstheorie.
\newblock \emph{Annalen der Physik}, 354(7):769--822, 1916.

\bibitem{Smyth1945}
Henry DeWolf Smyth.
\newblock \emph{Atomic Energy for Military Purposes}.
\newblock Princeton University Press, 1945.
\newblock \url{https://www.osti.gov/opennet/manhattan-project-history/publications/smyth_report.pdf}.

\bibitem{Franck1945}
James Franck et al.
\newblock The Franck Report: A Report to the Secretary of War.
\newblock June 1945.
\newblock \url{https://sgp.fas.org/eprint/franck.html}.

\bibitem{PlatoApology}
Plato.
\newblock \emph{Apology}.
\newblock Translated by Benjamin Jowett.
\newblock Project Gutenberg.
\newblock \url{https://www.gutenberg.org/ebooks/1656}.

\bibitem{Nakamoto2008}
Satoshi Nakamoto.
\newblock Bitcoin: A Peer-to-Peer Electronic Cash System.
\newblock 2008.
\newblock \url{https://bitcoin.org/bitcoin.pdf}.

\bibitem{Nakamoto2010WikiLeaks}
Satoshi Nakamoto.
\newblock Re: Wikileaks contact info?
\newblock BitcoinTalk, December 5, 2010.
\newblock Archived by the Satoshi Nakamoto Institute:
\url{https://satoshi.nakamotoinstitute.org/posts/bitcointalk/523/}.

\bibitem{Nakamoto2010DoS}
Satoshi Nakamoto.
\newblock Added some DoS limits, removed safe mode (0.3.19).
\newblock BitcoinTalk, December 12, 2010.
\newblock Archived by the Satoshi Nakamoto Institute:
\url{https://satoshi.nakamotoinstitute.org/posts/bitcointalk/543/}.

\bibitem{BitMEX2018}
BitMEX Research.
\newblock Does Satoshi have a million bitcoin?
\newblock August 20, 2018.
\newblock \url{https://www.bitmex.com/blog/satoshis-1-million-bitcoin}.

\end{thebibliography}

\end{document}
